\chapter{Counterterms at one loop and the quantum master equation at the UV boundary}

The one-loop counterterm, the analogue on $\mathbb R^4$ at order $\hbar$ of the counterterm of [C07, Thm B] (Remark \ref{rem:c4-thmB}), is fixed by the requirement that the order-$\hbar$ part of the flow $W(P(\varepsilon,\infty),I-I^{CT}[\varepsilon])$ have a limit as $\varepsilon\to0$ at generic leg momenta, and it is unique among the translation-invariant and $\mathfrak g$-invariant local functionals with coefficients in the span of $\varepsilon^{-1}$ and $\log\varepsilon$ (Propositions \ref{prop:c4-limit} and \ref{prop:c4-unique}). The regulator, however, cuts every edge at the Schwinger time $\varepsilon$ and breaks the gauge symmetry, so the one-loop counterterm contains the mass terms $\varepsilon^{-1}\int\mathrm{dvol}\,(A,A)$, whose $Q$-variation is nonzero. Transported to the UV boundary by Theorem \ref{thm:c3-uv}, the quantum master equation then fails by a finite local functional, and a finite local counterterm has to be chosen to remove it: this is the first renormalisation condition of the theory. This chapter computes $I^{CT}[\varepsilon]$ in every one-loop sector of superficial degree $\ge0$, identifies its $\log\varepsilon$ part as a rescaling of the fields with vanishing component along the coupling, and evaluates the master equation at the UV boundary in the sector with legs $(c,A|_V)$, where the first finite counterterm appears.

\section{The one-loop counterterm and its boundary form}

Throughout this chapter $\mathfrak g$ is simple. Then the Killing form $\operatorname{tr}_{\mathfrak g}(\mathrm{ad}_X\mathrm{ad}_Y)$ is a multiple of the invariant pairing, and the ratio $\operatorname{tr}_{\mathfrak g}(\mathrm{ad}_X\mathrm{ad}_Y)/(X,Y)$, for $X,Y\in\mathfrak g$ with $(X,Y)\neq0$, is a number independent of $X$ and $Y$. Every one-loop coefficient of this chapter is a rational multiple of
\[
\frac{\hbar\,\operatorname{tr}_{\mathfrak g}(\mathrm{ad}_X\mathrm{ad}_Y)}{16\pi^2\,(X,Y)},
\]
except the coefficients of the fourth-order adjoint traces of the quartic graphs, which are rational multiples of $\hbar/16\pi^2$. The upper Schwinger limit is $\infty$ (Chapter 3).

The value of a translation-invariant functional with $k$ legs on plane waves $e^{i(p_j,x)}$, $j=1,\dots,k$, with fibre elements as in Definition \ref{def:c1-vertex}, is $(2\pi)^4\delta(p_1+\dots+p_k)$ times a function of $(p_1,\dots,p_k)$ on the subspace where $p_1+\dots+p_k=0$; this function is its plane-wave value. Leg momenta are generic when every leg momentum is nonzero and no sum of the momenta of a nonempty proper subset of the legs vanishes. For $k=1$ the set of generic configurations is empty; for $k\ge2$ it is the complement of finitely many proper linear subspaces of the subspace $p_1+\dots+p_k=0$, an open dense subset. A functional depending on $\varepsilon$ has a limit as $\varepsilon\to0$ at generic leg momenta when the plane-wave values of its Taylor components converge as $\varepsilon\to0$ at every generic configuration of the leg momenta. Every limit of a flow in this chapter is a limit in this sense.

\begin{definition}[One-loop counterterm]
\label{def:c4-ict}
For a connected one-particle-irreducible graph $\gamma$ with legs, $b_1(\gamma)=1$ and all vertices labelled by $I$ (Proposition \ref{prop:c3-graphs}), let $w_\gamma(\varepsilon)$ be its value with every edge labelled by $P(\varepsilon,\infty)$ [C07, §15.4]. Every vertex of $\gamma$ lies on its cycle and carries a leg, so at generic leg momenta the momenta of distinct edges of $\gamma$ differ, and by Lemma \ref{lem:c4-local}, as $\varepsilon\to0$, the plane-wave value of $w_\gamma(\varepsilon)$ is the sum of multiples of negative integer powers of $\varepsilon$, a multiple of $\log\varepsilon$, and a function of $\varepsilon$ with a limit; the first two are polynomials in the leg momenta, and $\mathrm{Sing}_\varepsilon w_\gamma(\varepsilon)$ denotes the local functional whose plane-wave value is their sum [CG2, §13.7.2]. The one-loop counterterm is
\begin{equation}
I^{CT}[\varepsilon]=\hbar\sum_\gamma\frac{1}{|\operatorname{Aut}\gamma|}\,\mathrm{Sing}_\varepsilon\,w_\gamma(\varepsilon),
\label{eq:c4-ict}
\end{equation}
the sum over the connected one-particle-irreducible graphs $\gamma$ with $b_1(\gamma)=1$ and at least two legs, modulo constants. Its coefficients lie in the span of $\varepsilon^{-1}$ and $\log\varepsilon$ (Proposition \ref{prop:c4-structure}), the one-loop part of the renormalisation scheme $\mathscr A_{<0}$ [C07, Def 1.4.1; CG2, §13.7.2].
\end{definition}

The $\varepsilon$-dependence of a graph is carried by the lower ends of its Schwinger integrals. By \eqref{eq:c3-dP}, $\frac{d}{d\varepsilon}P(\varepsilon,\infty)=-L_\varepsilon$ with $L_\varepsilon=(Q^{GF}\otimes1)K_\varepsilon$.

\begin{proposition}[Boundary form of the scale derivative]
\label{prop:c4-boundary}
Let $\gamma$ be a graph with edge set $E(\gamma)$ [C07, §15.5], and let $w_\gamma(\varepsilon)$ be its value with every edge labelled by $P(\varepsilon,\infty)$. Then
\begin{equation}
\varepsilon\frac{d}{d\varepsilon}w_\gamma(\varepsilon)=\sum_{e\in E(\gamma)}w_\gamma\big(\text{edge }e\text{ labelled by }-\varepsilon L_\varepsilon,\ \text{the other edges by }P(\varepsilon,\infty)\big).
\label{eq:c4-boundary}
\end{equation}
In the Schwinger form $w_\gamma(\varepsilon)=\int_{[\varepsilon,\infty)^{E(\gamma)}}f_\gamma\prod_{e}dt_e$ of [C07, §15.5], with the edge $e$ carrying $L_{t_e}$, the $e$-th term of \eqref{eq:c4-boundary} is
\[
-\varepsilon\int_{[\varepsilon,\infty)^{E(\gamma)\setminus\{e\}}}f_\gamma\big|_{t_e=\varepsilon}\prod_{e'\neq e}dt_{e'} .
\]
\end{proposition}

\begin{proof}
The value $w_\gamma$ is multilinear in the kernels labelling the edges, so the Leibniz rule gives
\[
\frac{d}{d\varepsilon}w_\gamma=\sum_e w_\gamma\Big(\text{edge }e\text{ labelled by }\frac{d}{d\varepsilon}P(\varepsilon,\infty)\Big),
\]
and $\frac{d}{d\varepsilon}P(\varepsilon,\infty)=-L_\varepsilon$ by \eqref{eq:c3-dP}. In the Schwinger form, $\frac{d}{d\varepsilon}\int_\varepsilon^\infty(\cdots)\,dt_e=-(\cdots)|_{t_e=\varepsilon}$ for each edge.
\end{proof}

\begin{corollary}[Singular part from the boundary]
\label{cor:c4-sing}
At fixed leg momenta, if the right-hand side of \eqref{eq:c4-boundary} equals $-a\,\varepsilon^{-1}+b$ plus a term bounded by a multiple of $\varepsilon$, with $a$, $b$ independent of $\varepsilon$, then
\[
w_\gamma(\varepsilon)=a\,\varepsilon^{-1}+b\log\varepsilon+(\text{a function of }\varepsilon\text{ with a limit}),\qquad \mathrm{Sing}_\varepsilon w_\gamma(\varepsilon)=a\,\varepsilon^{-1}+b\log\varepsilon .
\]
\end{corollary}

\begin{proof}
$\frac{d}{d\varepsilon}\big(w_\gamma-a\varepsilon^{-1}-b\log\varepsilon\big)=\varepsilon^{-1}\big(\varepsilon\frac{d}{d\varepsilon}w_\gamma+a\varepsilon^{-1}-b\big)$ is $\varepsilon^{-1}$ times a term bounded by a multiple of $\varepsilon$, hence bounded near $\varepsilon=0$, so $w_\gamma-a\varepsilon^{-1}-b\log\varepsilon$ has a limit and differs from it by a term bounded by a multiple of $\varepsilon$.
\end{proof}

\fig{diagrams/composite/boundary_Leps_representatives.png}{Boundary graphs. (a), (b): the two terms of \eqref{eq:c4-boundary} for the bubble with two vertices $A|_V(A|_{W_{\mathbb R}^\perp})^2$ and external $A|_{W_{\mathbb R}^\perp},A|_{W_{\mathbb R}^\perp}$; the shaded edge is pinned at $t=\varepsilon$ and carries $L_\varepsilon=(Q^{GF}\otimes1)K_\varepsilon$. (c)--(e): the same mechanism for a triangle with one insertion of a composite operator.}

The singular coefficients of every graph are therefore computed from graphs with one edge pinned at $t=\varepsilon$. For one-loop graphs with at most two edges, the Schwinger integrals reduce to four elementary ones after a Gaussian integral over the loop momentum.

\begin{lemma}[Gaussian integral over the loop momentum]
\label{lem:c4-gauss}
Let $t_1,t_2>0$, $p\in V$, and let the numerator be a polynomial function of $p_1\in V$ of degree at most two. Then
\[
\begin{aligned}
&\int\frac{d^4p_1}{(2\pi)^4}\,e^{-t_1|p_1|^2-t_2|p_1-p|^2}\,(\text{numerator})\\
&\qquad=\frac{e^{-t_1t_2|p|^2/(t_1+t_2)}}{16\pi^2(t_1+t_2)^2}\Big[\text{numerator}-\frac{1}{4(t_1+t_2)}\big(\text{positive Laplacian of the numerator}\big)\Big]_{p_1=\frac{t_2}{t_1+t_2}p},
\end{aligned}
\]
the positive Laplacian, minus the sum of the second derivatives along an orthonormal basis, being taken in the variable $p_1$.
For one edge, $\int\frac{d^4p_1}{(2\pi)^4}e^{-t|p_1|^2}=(16\pi^2t^2)^{-1}$.
\end{lemma}

\begin{proof}
Completing the square and integrating the Gaussian,
\[
\begin{aligned}
&t_1|p_1|^2+t_2|p_1-p|^2=(t_1+t_2)\Big|p_1-\tfrac{t_2}{t_1+t_2}p\Big|^2+\tfrac{t_1t_2}{t_1+t_2}|p|^2,\\
&\int\frac{d^4p_1}{(2\pi)^4}e^{-(t_1+t_2)|p_1|^2}=\frac{\pi^2}{(2\pi)^4(t_1+t_2)^2}=\frac1{16\pi^2(t_1+t_2)^2}.
\end{aligned}
\] The normalised Gaussian average in the shifted variable has vanishing odd moments, and its second moment sends $(v,p_1)(v',p_1)$ to $(v,v')/(2(t_1+t_2))$ for $v,v'\in V$. A polynomial of degree at most two is its value at the centre plus a linear term plus its quadratic part; the average of the quadratic part is $\frac{1}{4(t_1+t_2)}$ times the sum of the second derivatives along an orthonormal basis, i.e. $-\frac{1}{4(t_1+t_2)}$ times the positive Laplacian of the numerator in $p_1$.
\end{proof}

\begin{lemma}[Elementary Schwinger integrals]
\label{lem:c4-integrals}
For $p\neq0$, modulo functions of $\varepsilon$ with a limit as $\varepsilon\to0$ ($\mathscr A_{\ge0}$ of [C07, Def 1.4.1]),
\[
\begin{aligned}
&\int_\varepsilon^\infty\frac{dt}{t^2}=\varepsilon^{-1},\\
&\iint_{[\varepsilon,\infty)^2}\frac{e^{-t_1t_2|p|^2/(t_1+t_2)}}{(t_1+t_2)^3}\,dt_1dt_2=\tfrac14\varepsilon^{-1}+\tfrac16|p|^2\log\varepsilon,\\
&\iint_{[\varepsilon,\infty)^2}\frac{e^{-t_1t_2|p|^2/(t_1+t_2)}}{(t_1+t_2)^2}\,dt_1dt_2=-\log\varepsilon,\\
&\iint_{[\varepsilon,\infty)^2}\frac{t_1t_2\,e^{-t_1t_2|p|^2/(t_1+t_2)}}{(t_1+t_2)^4}\,dt_1dt_2=-\tfrac16\log\varepsilon.
\end{aligned}
\]
\end{lemma}

\begin{proof}
The integrands are symmetric in $t_1,t_2$, so by Proposition \ref{prop:c4-boundary} each $\varepsilon\frac{d}{d\varepsilon}$ equals $-2\varepsilon$ times the integral over $t\in[\varepsilon,\infty)$ of the integrand at $t_1=\varepsilon$, $t_2=t$. On the pinned edge the exponent satisfies $0\le\frac{\varepsilon t|p|^2}{\varepsilon+t}\le\varepsilon|p|^2$, so $e^{-\varepsilon t|p|^2/(\varepsilon+t)}$ equals $1-\frac{\varepsilon t|p|^2}{\varepsilon+t}$ plus a remainder between $0$ and $\tfrac12\varepsilon^2|p|^4$, uniformly in $t$. With
\[
\begin{aligned}
&\int_\varepsilon^\infty\frac{dt}{(\varepsilon+t)^2}=\frac1{2\varepsilon},\qquad \int_\varepsilon^\infty\frac{dt}{(\varepsilon+t)^3}=\frac1{8\varepsilon^2},\\
&\int_\varepsilon^\infty\frac{t\,dt}{(\varepsilon+t)^4}=\int_\varepsilon^\infty\Big(\frac1{(\varepsilon+t)^3}-\frac{\varepsilon}{(\varepsilon+t)^4}\Big)dt=\frac1{8\varepsilon^2}-\frac1{24\varepsilon^2}=\frac1{12\varepsilon^2},
\end{aligned}
\]
the three pinned integrals are, up to terms bounded by a multiple of $\varepsilon$,
\[
\begin{aligned}
2\varepsilon\int_\varepsilon^\infty\frac{e^{-\varepsilon t|p|^2/(\varepsilon+t)}}{(\varepsilon+t)^3}dt&=\frac{1}{4\varepsilon}-\frac{|p|^2}{6},\\
2\varepsilon\int_\varepsilon^\infty\frac{e^{-\varepsilon t|p|^2/(\varepsilon+t)}}{(\varepsilon+t)^2}dt&=1,\qquad
2\varepsilon\int_\varepsilon^\infty\frac{\varepsilon t\,e^{-\varepsilon t|p|^2/(\varepsilon+t)}}{(\varepsilon+t)^4}dt=\frac16 .
\end{aligned}
\]
Corollary \ref{cor:c4-sing} with $(a,b)$ equal to $(\tfrac14,\tfrac16|p|^2)$, $(0,-1)$, $(0,-\tfrac16)$ gives the three expansions. The one-edge integral is exact.
\end{proof}

\begin{lemma}[Singular part of a one-loop graph]
\label{lem:c4-local}
Let $\gamma$ be a connected one-particle-irreducible graph with $b_1(\gamma)=1$, so that every vertex lies on the cycle and $|E(\gamma)|=|V(\gamma)|$. Let its vertex functions be polynomials in the momenta, up to a factor depending on the leg momenta alone (such as the Fourier transform of a four-form at minus the total leg momentum, Chapter 5), and let its numerator, the product of the vertex functions and of the polynomial factors of the edge components of $P(\varepsilon,\infty)$ (Proposition \ref{prop:c3-P}), be homogeneous in the loop momentum and the leg momenta; for vertices labelled by $I$ and legs in fixed slots this holds by Lemma \ref{lem:c4-weights}. Its superficial degree is the degree of the numerator plus $4-2|E(\gamma)|$. Let the leg momenta be such that the momenta of any two distinct edges differ by a nonzero vector. Then, at these leg momenta and as $\varepsilon\to0$,
\[
w_\gamma(\varepsilon)=\sum_{1\le k\le\frac12(\text{superficial degree})}\varepsilon^{-k}\,a_k+b\,\log\varepsilon+(\text{a function of }\varepsilon\text{ with a limit}),
\]
where $a_k$ and $b$ are polynomials in the leg momenta, $a_k$ homogeneous of degree (superficial degree)$-2k$ and $b$ homogeneous of degree equal to the superficial degree, and $\varepsilon\frac{d}{d\varepsilon}$ of the function with a limit tends to $0$. In particular the sum of the terms with $\varepsilon^{-k}$ and $\log\varepsilon$ is the plane-wave value of a local functional, $\mathrm{Sing}_\varepsilon w_\gamma(\varepsilon)$, and it vanishes when the superficial degree is negative.
\end{lemma}

\begin{proof}
Let the edge $e\in E(\gamma)$ carry the momentum $p_1-q_e$, with $p_1$ the loop momentum and $q_e$ a sum of leg momenta, so that $q_e\neq q_f$ for $e\neq f$. By Proposition \ref{prop:c3-P}, $w_\gamma(\varepsilon)$ is the integral over $p_1$, with the measure $d^4p_1/(2\pi)^4$, of the numerator times $\prod_ee^{-\varepsilon|p_1-q_e|^2}/|p_1-q_e|^2$, and $e^{-\varepsilon|k|^2}/|k|^2=\int_\varepsilon^\infty e^{-t|k|^2}dt$ introduces a Schwinger time $t_e$ for each edge. By Tonelli's theorem,
\[
\int\frac{d^4p_1}{(2\pi)^4}\int_{[\varepsilon,\infty)^{E(\gamma)}}\big|\text{numerator}\big|\,e^{-\sum_et_e|p_1-q_e|^2}\prod_edt_e=\int\frac{d^4p_1}{(2\pi)^4}\,\big|\text{numerator}\big|\prod_e\frac{e^{-\varepsilon|p_1-q_e|^2}}{|p_1-q_e|^2},
\]
and the right side is finite: the points $p_1=q_e$ are distinct, near each of them exactly one factor is singular, $|p_1-q_e|^{-2}$ being integrable near $p_1=q_e$ in four dimensions, and the Gaussian factors bound the polynomial numerator at large $|p_1|$. By Fubini's theorem, $w_\gamma(\varepsilon)$ is the integral over $t\in[\varepsilon,\infty)^{E(\gamma)}$ of the integral over $p_1$, with the measure $\prod_edt_e\,d^4p_1/(2\pi)^4$, of $\exp\big(-\sum_et_e|p_1-q_e|^2\big)$ times the numerator. Write $\tau=\sum_et_e$. Completing the square,
\[
\sum_et_e|p_1-q_e|^2=\tau\Big|p_1-\sum_e\frac{t_e}{\tau}\,q_e\Big|^2+\tau\cdot\frac12\sum_{e,f}\frac{t_e}{\tau}\frac{t_f}{\tau}|q_e-q_f|^2 ,
\]
and the Gaussian integral over $p_1$, with the moments of Lemma \ref{lem:c4-gauss} (each contracted pair of loop momenta gives $(2\tau)^{-1}$ times the metric), gives
\[
w_\gamma(\varepsilon)=\sum_{j\ge0}\int_{[\varepsilon,\infty)^{E(\gamma)}}\frac{(\text{numerator})_j}{16\pi^2\,\tau^{2+j}}\,\exp\Big(-\tau\cdot\frac12\sum_{e,f}\frac{t_e}{\tau}\frac{t_f}{\tau}|q_e-q_f|^2\Big)\prod_edt_e ,
\]
where $(\text{numerator})_j$, the part with $j$ contracted pairs, is a polynomial in the ratios $t_e/\tau$ and the leg momenta, homogeneous of degree (degree of the numerator)$-2j$, and zero when $2j$ exceeds the degree of the numerator. Put $m_j=j+2-|E(\gamma)|$; then $(\text{numerator})_j\neq0$ implies $2m_j\le$ (superficial degree). The ratios $t_e/\tau$ range over the simplex $\sum_et_e/\tau=1$, and $\prod_edt_e$ is $\tau^{|E(\gamma)|-1}d\tau$ times the Lebesgue measure of the simplex. Since $\sum_{e\neq f}t_et_f/\tau^2=1-\sum_et_e^2/\tau^2\ge1-\max_et_e/\tau$, the factor multiplying $\tau$ in the exponent is at least one half of the smallest of the numbers $|q_e-q_f|^2$, $e\neq f$, times $1-\max_et_e/\tau$, and $\tau(1-\max_et_e/\tau)$ is the sum of the Schwinger times other than the largest.

The region $\tau\ge1$. The integrand of the $j$-th term is bounded by a multiple of $\tau^{-2-j}$ times the exponential of minus a positive multiple of the sum of the Schwinger times other than the largest. Integrating these times against the exponential, and the largest time, which is at least $\tau/|E(\gamma)|\ge1/|E(\gamma)|$, against $\tau^{-2-j}$, gives a bound independent of $\varepsilon$; by dominated convergence this part has a limit as $\varepsilon\to0$.

The region $\tau\le1$. If $m_j<0$, the integrand times $\tau^{|E(\gamma)|-1}$ is bounded by a multiple of $\tau^{-1-m_j}$, which is integrable, and the part has a limit. If $m_j\ge0$, expand the exponential to order $m_j$. The remainder is bounded by $1/(m_j+1)!$ times the $(m_j+1)$-th power of the exponent, so its contribution is bounded by a multiple of $\tau^{-2-j+m_j+1}\tau^{|E(\gamma)|-1}=\tau^0$ and has a limit. The term of order $i\le m_j$ is the integral, over the part of the simplex where $\min_et_e/\tau\ge\varepsilon$, of $(16\pi^2\,i!)^{-1}(\text{numerator})_j$ times the $i$-th power of minus the factor multiplying $\tau$ in the exponent, times
\[
\int_{\varepsilon/\min_e(t_e/\tau)}^{1}\tau^{-1-(m_j-i)}\,d\tau=\begin{cases}\dfrac{\varepsilon^{-(m_j-i)}\big(\min_et_e/\tau\big)^{m_j-i}-1}{m_j-i}, & i<m_j,\\[2mm] \log\big(\min_et_e/\tau\big)-\log\varepsilon, & i=m_j.\end{cases}
\]
The part of the simplex where $\min_et_e/\tau<\varepsilon$ has measure at most a multiple of $\varepsilon$, the function $(\min_et_e/\tau)^{m_j-i}$ is at most $\varepsilon^{m_j-i}$ there, and $\log(\min_et_e/\tau)$ is integrable on the simplex. Hence the term of order $i<m_j$ is $\varepsilon^{-(m_j-i)}$ times an integral over the whole simplex plus a function with a limit, and the term of order $m_j$ is $-\log\varepsilon$ times an integral over the whole simplex plus a function with a limit. These integrals over the simplex, of polynomials in the ratios and the leg momenta times $(\min_et_e/\tau)^{m_j-i}$, are polynomials in the leg momenta, homogeneous of degree (degree of the numerator)$-2j+2i$, which equals (superficial degree)$-2(m_j-i)$. Collecting the powers $k=m_j-i$ gives the expansion, with $1\le k\le m_j\le\frac12$(superficial degree); for negative superficial degree every $m_j$ with $(\text{numerator})_j\neq0$ is negative, and the singular part vanishes.

Finally, $\varepsilon\frac{d}{d\varepsilon}$ of the function with a limit tends to $0$. That function is the sum, over $j$ and $i$, of the following pieces; in the second and third item the simplex integrand of the term of order $i$ is $(16\pi^2\,i!)^{-1}(\text{numerator})_j$ times the $i$-th power of minus the factor multiplying $\tau$ in the exponent, a polynomial in the ratios and the leg momenta.
\begin{itemize}
\item The integrals over $t\in[\varepsilon,\infty)^{E(\gamma)}$ restricted to $\tau\ge1$, those restricted to $\tau\le1$ with $m_j<0$, and the remainders of the expansion of the exponential on $\tau\le1$. By Proposition \ref{prop:c4-boundary}, $\varepsilon\frac{d}{d\varepsilon}$ of each is minus $\varepsilon$ times its integrals over the faces $t_e=\varepsilon$; by the bounds above these face integrals are bounded on $\tau\ge1$ and bounded by a multiple of $1+|\log\varepsilon|$ on $\tau\le1$, the integrand being at most a multiple of $\tau^{1-|E(\gamma)|}$ there.
\item For the term of order $i<m_j$, with $k=m_j-i$: $-k^{-1}\varepsilon^{-k}h(\varepsilon)$, where $h(\varepsilon)$ is the integral of the simplex integrand times $(\min_et_e/\tau)^k$ over the part of the simplex where $\min_et_e/\tau<\varepsilon$; and $-k^{-1}$ times the integral of the simplex integrand over the part where $\min_et_e/\tau\ge\varepsilon$, which comes from the $-1$ in the integral over $\tau$.
\item For the term of order $m_j$: $h(\varepsilon)\log\varepsilon$, where $h(\varepsilon)$ is the integral of the simplex integrand over the part where $\min_et_e/\tau<\varepsilon$; and the integral of the simplex integrand times $\log(\min_et_e/\tau)$ over the part where $\min_et_e/\tau\ge\varepsilon$.
\end{itemize}
The part of the simplex where $\min_et_e/\tau<s$ is the union over $e$ of the parts where $t_e/\tau<s$. For a function continuous on the part where $\min_et_e/\tau\ge s_0$, its integral over the part where $\min_et_e/\tau\ge s$, $s\ge s_0$, is therefore differentiable in $s$, with derivative minus a sum of integrals of the function over the faces where $\min_et_e/\tau=s$, each face having $(|E(\gamma)|-2)$-dimensional measure at most that of a face of the simplex; the same holds, with the opposite sign and $s_0=0$, for the integral of a function continuous on the simplex over the part where $\min_et_e/\tau<s$, and that integral is at most the maximum of the absolute value of the function there times a multiple of $s$. In the second item the integrand of $h$ is at most a multiple of $\varepsilon^k$ where $\min_et_e/\tau<\varepsilon$, so $h(\varepsilon)$ and $\varepsilon\frac{dh}{d\varepsilon}$ are bounded by multiples of $\varepsilon^{k+1}$, and $\varepsilon\frac{d}{d\varepsilon}\big(\varepsilon^{-k}h(\varepsilon)\big)=-k\varepsilon^{-k}h(\varepsilon)+\varepsilon^{1-k}\frac{dh}{d\varepsilon}$ is bounded by a multiple of $\varepsilon$; $\varepsilon\frac{d}{d\varepsilon}$ of the integral over the part where $\min_et_e/\tau\ge\varepsilon$ is bounded by a multiple of $\varepsilon$. In the third item $h(\varepsilon)$ and $\varepsilon\frac{dh}{d\varepsilon}$ are bounded by multiples of $\varepsilon$, so $\varepsilon\frac{d}{d\varepsilon}\big(h(\varepsilon)\log\varepsilon\big)=h(\varepsilon)+\varepsilon\frac{dh}{d\varepsilon}\log\varepsilon$ is bounded by a multiple of $\varepsilon(1+|\log\varepsilon|)$; and $\log(\min_et_e/\tau)=\log\varepsilon$ on the faces where $\min_et_e/\tau=\varepsilon$, so $\varepsilon\frac{d}{d\varepsilon}$ of the integral over the part where $\min_et_e/\tau\ge\varepsilon$ is bounded by a multiple of $\varepsilon|\log\varepsilon|$. Every piece has $\varepsilon\frac{d}{d\varepsilon}$ tending to $0$.
\end{proof}

\begin{lemma}[One-loop graphs as trees with a one-loop vertex]
\label{lem:c4-factor}
Consider the graphs of Proposition \ref{prop:c3-graphs}, with labelled legs, whose vertices are labelled by Taylor components of $I$, except at most one vertex of valence at least $1$, the marked vertex, labelled by a Taylor component of a further functional (in Chapter 5, the operator $O$). For a connected graph $\Gamma$ with $b_1(\Gamma)=1$, let $\gamma$ be its cycle together with the vertices on the cycle, the remaining half-edges at these vertices being the legs of $\gamma$, and let $T$ be the tree obtained from $\Gamma$ by contracting $\gamma$ to one vertex $s$, whose half-edges are the legs of $\gamma$.
\begin{itemize}
\item $\gamma$ is connected and one-particle-irreducible with $b_1(\gamma)=1$. The automorphisms of $T$ fixing the legs are the identity, and so are those of every tree with the given legs whose vertices are labelled as above except one distinguished vertex $s$. Restriction to $\gamma$ is a bijection from the automorphisms of $\Gamma$ fixing the legs to the automorphisms of $\gamma$ fixing the legs of $\gamma$. The value $w_\Gamma$ is $w_T$ with $s$ labelled by the multilinear form $w_\gamma$ in the half-edges of $s$.
\item The sum of $w_\Gamma/|\operatorname{Aut}\Gamma|$ over the connected graphs $\Gamma$ with $b_1(\Gamma)=1$ and the given legs equals the sum of $w_T/|\operatorname{Aut}T|=w_T$ over the trees $T$ with the given legs and one vertex $s$, labelled by the Taylor component, with legs the half-edges of $s$, of the sum of $w_\gamma/|\operatorname{Aut}\gamma|$ over the connected one-particle-irreducible graphs $\gamma$ with $b_1(\gamma)=1$.
\item A connected one-particle-irreducible graph with $b_1=1$, one leg and all vertices labelled by $I$ is a cubic vertex of $I$ with one edge from the vertex to itself, and its value is $0$ before any integration over positions or momenta.
\end{itemize}
\end{lemma}

\begin{proof}
First item. A connected graph with $b_1=1$ contains exactly one cycle. Every vertex of $\gamma$ lies on it, so removing an edge of $\gamma$ leaves $\gamma$ connected, and $b_1(\gamma)=|E(\gamma)|-|V(\gamma)|+1=1$. Contracting the cycle lowers $b_1$ by one, so $T$ is a tree.

Let $T$ be such a tree, with distinguished vertex $s$. An automorphism $\alpha$ of $T$ fixing the legs preserves the labels of the vertices, so it fixes $s$ and the marked vertex. For a vertex $v\neq s$, the vertices beyond $v$ are $v$ and the vertices separated from $s$ by $v$. A vertex beyond $v$ at maximal distance from $s$ has one half-edge towards $s$ and its other half-edges are legs; if it is labelled by $I$ it has at least two legs. Hence the vertices beyond $v$ carry a leg or include the marked vertex. Let $v^{\prime}\neq v$ be a second vertex different from $s$, the two being exchanged if $v$ lies beyond $v^{\prime}$. If $v^{\prime}$ lies beyond $v$ and $v$ is labelled by $I$, then $v$ has, besides its half-edges towards $s$ and towards $v^{\prime}$, a further half-edge, which is a leg or leads to vertices beyond $v$ and not beyond $v^{\prime}$ that carry a leg or include the marked vertex; if $v^{\prime}$ lies beyond $v$ and $v$ is the marked vertex, the vertices beyond $v$ include the marked vertex and those beyond $v^{\prime}$ exclude it; if neither lies beyond the other, the two sets of vertices beyond them are disjoint, and each carries a leg or includes the marked vertex. In every case $v$ and $v^{\prime}$ differ in the set of legs carried by the vertices beyond them or in whether these include the marked vertex. Both data are preserved by $\alpha$, so $\alpha$ fixes every vertex. A tree has no two edges with the same ends and no edge from a vertex to itself, so $\alpha$ fixes every edge, and, with the legs, every half-edge.

An automorphism of $\Gamma$ fixing the legs maps the cycle to itself and therefore induces an automorphism of $T$ fixing the legs, which is the identity. It fixes every vertex outside the cycle and every half-edge outside the edges of the cycle, the legs of $\gamma$ included, so it restricts to an automorphism of $\gamma$ fixing the legs of $\gamma$, and this restriction determines it. Conversely, an automorphism of $\gamma$ fixing its legs extends by the identity to an automorphism of $\Gamma$ fixing the legs.

The value $w_\Gamma$ is the contraction of the tensor product of the vertex functions along the edges of $\Gamma$, with the edge rule of Proposition \ref{prop:c3-graphs}. The contractions along different edges commute [C07, §6.4], so contracting the edges of the cycle first gives $w_\gamma$, a multilinear form in the half-edges of $s$, and the remaining contractions give $w_T$ with $s$ labelled by $w_\gamma$.

Second item. By the first item the half-edges of $s$ are determined by $T$ and its legs, and $\Gamma\mapsto(T,\gamma)$, with the legs of $\gamma$ identified with the half-edges of $s$, is a bijection from the isomorphism classes of connected graphs with $b_1=1$ and the given legs to the pairs consisting of an isomorphism class of trees $T$ with the given legs and a vertex $s$, and an isomorphism class of connected one-particle-irreducible graphs $\gamma$ with $b_1(\gamma)=1$ whose legs are the half-edges of $s$; the inverse glues $\gamma$ into $T$ in place of $s$. With $|\operatorname{Aut}\Gamma|=|\operatorname{Aut}\gamma|$, $|\operatorname{Aut}T|=1$ and the linearity of $w_T$ in the label of $s$,
\[
\sum_\Gamma\frac{w_\Gamma}{|\operatorname{Aut}\Gamma|}=\sum_T\sum_\gamma\frac{1}{|\operatorname{Aut}\gamma|}\,w_T\big(s\text{ labelled by }w_\gamma\big)=\sum_T\frac{1}{|\operatorname{Aut}T|}\,w_T\Big(s\text{ labelled by }\sum_\gamma\frac{w_\gamma}{|\operatorname{Aut}\gamma|}\Big).
\]

Third item. A cycle with $n$ vertices of valence at least three has at least $n$ legs, so one leg forces $n=1$ and a cubic vertex, two of whose slots are joined by the edge. By Definition \ref{def:c1-vertex}, every cubic vertex function of $I$, namely $A\,A\,A$ (which contains $(A|_V)^3$ and $A|_V(A|_{W_{\mathbb R}^\perp})^2$), $\lambda\,\lambda\,A$, $A\,c\,A^*$, $c\,\lambda\,\lambda^*$ and $c\,c\,c^*$, is the form $(X_1,[X_2,X_3])$ in the $\mathfrak g$-factors $X_1,X_2,X_3$ of its three legs, up to sign, times a factor independent of them; by the invariance of $(-,-)$ and the antisymmetry of the bracket, $(X_1,[X_2,X_3])$ is totally antisymmetric. The $\mathfrak g$-factor of every component of $P(\varepsilon,\infty)$ is a multiple of $(-,-)^{-1}$ (Proposition \ref{prop:c3-P}), which is symmetric. The contraction of two arguments of a totally antisymmetric form with a symmetric tensor is $0$, so the vertex with two of its slots joined is the zero functional, before any integration.
\end{proof}

\begin{proposition}[Existence of the limit]
\label{prop:c4-limit}
The order-$\hbar$ part of $W(P(\varepsilon,\infty),I-I^{CT}[\varepsilon])$ has a limit as $\varepsilon\to0$ at generic leg momenta.
\end{proposition}

\begin{proof}
By Proposition \ref{prop:c3-graphs} with $J=I-I^{CT}[\varepsilon]$, the order-$\hbar$ part is $\hbar$ times the sum of $w_\Gamma/|\operatorname{Aut}\Gamma|$ over the connected graphs $\Gamma$ with $b_1(\Gamma)=1$ and vertices labelled by $I$, plus $\hbar$ times the sum of $w_T/|\operatorname{Aut}T|$ over the trees $T$ with one vertex $s$ labelled by a Taylor component of $-\hbar^{-1}I^{CT}[\varepsilon]$ and the other vertices labelled by $I$. By Lemma \ref{lem:c4-factor} (without marked vertex), $|\operatorname{Aut}T|=1$ for these trees, and the first sum is the sum over the same trees of $w_T$ with $s$ labelled by the sum of $w_\gamma/|\operatorname{Aut}\gamma|$ over the connected one-particle-irreducible graphs $\gamma$ with $b_1(\gamma)=1$. A cycle of vertices of valence at least three carries legs, so every such $\gamma$ has legs; those with one leg have value $0$ (Lemma \ref{lem:c4-factor}), and $I^{CT}[\varepsilon]$ has no Taylor component with one leg (Definition \ref{def:c4-ict}). By \eqref{eq:c4-ict}, the order-$\hbar$ part is therefore $\hbar$ times the sum over the trees $T$ whose vertex $s$ has at least two half-edges of $w_T$ with $s$ labelled by
\[
\sum_\gamma\frac{1}{|\operatorname{Aut}\gamma|}\Big(w_\gamma(\varepsilon)-\mathrm{Sing}_\varepsilon w_\gamma(\varepsilon)\Big),
\]
the sum over the $\gamma$ whose legs are the half-edges of $s$. For a fixed number of legs this sum is finite: a tree with $r$ vertices labelled by $I$ and a vertex $s$ with $m$ half-edges has $r$ edges and at least $3r+m-2r=r+m$ legs, and a cycle with $m$ legs has at most $m$ vertices.

Fix generic leg momenta. Each edge of $T$ separates $T$ into two parts, and each part carries a leg: in the part without $s$, a vertex at maximal distance from $s$ carries at least two legs; in the part with $s$, the vertex $s$ has, besides the half-edge on the path from $s$ to the edge, a further half-edge, which is a leg or leads to vertices separated from $s$ and from the edge, among which a vertex at maximal distance from $s$ carries at least two legs. Hence the edge carries the sum of the momenta of a nonempty proper subset of the legs, a nonzero momentum $k$, and its edge factor, a polynomial times $e^{-\varepsilon|k|^2}/|k|^2$ (Proposition \ref{prop:c3-P}), converges to the polynomial times $|k|^{-2}$. The half-edges of $s$ carry the sums of the momenta of disjoint nonempty sets of legs, one set for each half-edge. Two distinct edges of the cycle of $\gamma$ cut it into two arcs, each containing a vertex, and every vertex of $\gamma$ carries a leg of $\gamma$; the momenta of the two edges differ by the sum of the momenta entering $\gamma$ along one arc, the sum of the momenta of a nonempty proper subset of the legs of $T$, which is nonzero. By Lemma \ref{lem:c4-local}, with the homogeneity of Lemma \ref{lem:c4-weights}, the plane-wave value of $w_\gamma(\varepsilon)-\mathrm{Sing}_\varepsilon w_\gamma(\varepsilon)$ at these momenta has a limit as $\varepsilon\to0$. The plane-wave value of $w_T$ is the sum, over the components of the internal slots, of the products of the vertex functions of $T$, of its edge factors and of this value, a finite sum of products of functions of $\varepsilon$ with limits.
\end{proof}

\begin{remark}[Relation with Theorem B]
\label{rem:c4-thmB}
[C07, Thm B] constructs, on a compact manifold with a finite upper scale, a counterterm at all orders in $\hbar$, unique among the local functionals with coefficients in $\mathscr A_{<0}$, for which the flow of the interaction minus the counterterm has a limit as $\varepsilon\to0$ as a functional; at one loop, its part with $k$ legs is the singular part of the order-$\hbar$, $k$-leg part of the flow of the interaction minus the counterterms with fewer legs, and it is independent of the upper scale [C07, §1.5]. Definition \ref{def:c4-ict} and Propositions \ref{prop:c4-limit} and \ref{prop:c4-unique} give its analogue on $\mathbb R^4$ with upper scale $\infty$, for the even, local, at least cubic interaction $I$ of \eqref{eq:c1-I}, in the following scope: the order is $\hbar$; the singular parts are computed graph by graph on the one-particle-irreducible graphs, and the one-particle-reducible graphs carry them through their cycles (Lemma \ref{lem:c4-factor}); the limit is the limit of plane-wave values at generic leg momenta; and uniqueness holds among the translation-invariant and $\mathfrak g$-invariant local functionals whose coefficients lie in a space of functions of $\varepsilon$ containing $\varepsilon^{-1}$ and $\log\varepsilon$ in which every function with a limit as $\varepsilon\to0$ is zero.
\end{remark}

\begin{proposition}[Uniqueness of the one-loop counterterm]
\label{prop:c4-unique}
Let $I^{\prime}$ be a local functional of order $\hbar$, without constant term, invariant under translations and under the adjoint action of $\mathfrak g$, whose coefficients lie in a space of functions of $\varepsilon$ containing $\varepsilon^{-1}$ and $\log\varepsilon$ in which every function with a limit as $\varepsilon\to0$ is zero, for instance the span of $\varepsilon^{-1}$ and $\log\varepsilon$, or a complement $\mathscr A_{<0}$ of $\mathscr A_{\ge0}$ containing them [C07, Def 1.4.1]. If the order-$\hbar$ part of $W(P(\varepsilon,\infty),I-I^{\prime})$ has a limit as $\varepsilon\to0$ at generic leg momenta, then $I^{\prime}=I^{CT}[\varepsilon]$.
\end{proposition}

\begin{proof}
The vertex functions of $I$ and the components of $P(\varepsilon,\infty)$ are translation invariant, and their $\mathfrak g$-factors are built from the bracket, $(-,-)$ and $(-,-)^{-1}$; so $I^{CT}[\varepsilon]$, and with it the difference $I^{CT}[\varepsilon]-I^{\prime}$, is invariant under translations and under the adjoint action of $\mathfrak g$, with coefficients in the given space of functions. A translation-invariant local functional with one leg is, modulo total derivatives, the integral of a constant linear form on the fibre of the bundle of fields. The fibre is $\mathfrak g$ tensored with a representation of $\mathrm{Spin}(4,\mathbb C)\times\mathrm{Spin}(6,\mathbb C)$ (Chapter 1), so a $\mathfrak g$-invariant linear form on it lies in the tensor product of the invariant linear forms on $\mathfrak g$ with the dual of that representation; the invariant linear forms on $\mathfrak g$ are, through $(-,-)$, the centre of $\mathfrak g$, which is $0$ for $\mathfrak g$ simple. Hence the difference has no Taylor component with one leg, and by hypothesis none without legs.

The order-$\hbar$ parts of $W(P(\varepsilon,\infty),I-I^{\prime})$ and $W(P(\varepsilon,\infty),I-I^{CT}[\varepsilon])$ differ by $\hbar$ times the sum of the trees with one vertex labelled by a Taylor component of $\hbar^{-1}(I^{CT}[\varepsilon]-I^{\prime})$ and all other vertices labelled by $I$ (Proposition \ref{prop:c3-graphs}). That vertex has at least two half-edges, so, as in the proof of Proposition \ref{prop:c4-limit}, every edge of such a tree carries a nonzero momentum at generic leg momenta. By the hypothesis and Proposition \ref{prop:c4-limit}, the sum of the trees has a limit at generic leg momenta. Suppose the difference is nonzero, and let $k\ge2$ be the smallest number of legs of its nonzero Taylor components. A tree with $r\ge1$ vertices labelled by $I$, each with at least three half-edges, has $r$ edges, and its vertex labelled by the difference has at least $k$ half-edges, so it has at least $k+3r-2r=k+r$ legs; the $k$-leg part of the sum of the trees is therefore the $k$-leg Taylor component of $\hbar^{-1}(I^{CT}[\varepsilon]-I^{\prime})$ itself, which then has a limit at generic leg momenta. Its plane-wave value is a finite sum of products of linearly independent coefficient functions of $\varepsilon$ with polynomials in $(p_1,\dots,p_k)$ on the subspace where $p_1+\dots+p_k=0$. At a generic configuration this sum is a function of $\varepsilon$ in the given space with a limit, hence zero, and by linear independence every polynomial vanishes there. The generic configurations are dense in the subspace, so the polynomials vanish on it. A polynomial vanishing on that subspace is a combination of the components of $p_1+\dots+p_k$ with polynomial coefficients, the plane-wave value of a total derivative; so the $k$-leg component is zero, a contradiction.
\end{proof}

\section{Two-point graphs}

A quadratic functional is evaluated on plane waves. For two legs of $A|_V$ the background is $A(p)e^{i(p,x)}+A(-p)e^{-i(p,x)}$ with $A(\pm p)\in V^\vee\otimes\mathfrak g$, and the plane-wave value of a functional is the coefficient of $(2\pi)^4\delta(0)$ in its part bilinear in $A(p)$ and $A(-p)$; for two odd legs the background is $r_1(\cdot)e^{i(p,x)}+r_2(\cdot)e^{-i(p,x)}$ and the plane-wave value is the coefficient of $r_1r_2$, as in Definition \ref{def:c1-vertex}. Since $d$ becomes $i(p,-)$ and $d^*$ becomes $-i\,\iota_p$,
\begin{equation}
\begin{aligned}
\int\mathrm{dvol}\,(A|_V,A|_V)&\mapsto2\,(A(p),A(-p)),\qquad \int\mathrm{dvol}\,(d^*A|_V,d^*A|_V)\mapsto2\,(\iota_pA(p),\iota_pA(-p)),\\
\int\mathrm{dvol}\,(dA|_V,dA|_V)&\mapsto2\,\big(|p|^2(A(p),A(-p))-(\iota_pA(p),\iota_pA(-p))\big),\\
\int\mathrm{dvol}\,(dc,A^*)&\mapsto i\,(p\,c(-p),A^*(p))\quad\text{for the legs }r_1A^*(p)e^{i(p,x)}+r_2c(-p)e^{-i(p,x)},
\end{aligned}
\label{eq:c4-pw}
\end{equation}
where $(A(p),A(-p))$ pairs the $V_{10}^\vee$ factors by the metric and the $\mathfrak g$ factors by $(-,-)$, $(p\,c(-p),A^*(p))$ is the fibre pairing of $p\otimes c(-p)\in V\otimes\mathfrak g$ with $A^*(p)$, and the Hodge normalisation of Definition \ref{def:c1-action} gives $((p,-)\wedge A(p),(p,-)\wedge A(-p))=|p|^2(A(p),A(-p))-(\iota_pA(p),\iota_pA(-p))$.

\fig{diagrams/composite/twopoint_AV.png}{The one-loop one-particle-irreducible graphs with external $A|_V,A|_V$; every edge carries $P(\varepsilon,\infty)$. (a), (b): quartic vertex with an $A|_V$, resp. $A|_{W_{\mathbb R}^\perp}$, loop. (c): gauge bubble with $A|_V$ lines. (d): bubble with $A|_{W_{\mathbb R}^\perp}$ lines. (e): gaugino loop. (f): ghost loop of $c$--$A^*$ lines.}

\begin{tabular}{l|c|c|c}
graph & $1/|\operatorname{Aut}\gamma|$ & $\varepsilon^{-1}$ part & $\log\varepsilon$ part \\
\hline
(a) & $\tfrac12$ & $\tfrac32\int\mathrm{dvol}\,(A|_V,A|_V)$ & $0$ \\
(b) & $\tfrac12$ & $3\int\mathrm{dvol}\,(A|_V,A|_V)$ & $0$ \\
(c) & $\tfrac12$ & $-\tfrac{9}{16}\int\mathrm{dvol}\,(A|_V,A|_V)$ & $\tfrac{19}{24}\int\mathrm{dvol}\,(dA|_V,dA|_V)-\tfrac18\int\mathrm{dvol}\,(d^*A|_V,d^*A|_V)$ \\
(d) & $\tfrac12$ & $-\tfrac34\int\mathrm{dvol}\,(A|_V,A|_V)$ & $-\tfrac12\int\mathrm{dvol}\,(dA|_V,dA|_V)$ \\
(e) & $\tfrac12$ & $-\int\mathrm{dvol}\,(A|_V,A|_V)$ & $-\tfrac43\int\mathrm{dvol}\,(dA|_V,dA|_V)$ \\
(f) & $1$ & $\tfrac{1}{16}\int\mathrm{dvol}\,(A|_V,A|_V)$ & $\tfrac{1}{24}\int\mathrm{dvol}\,(dA|_V,dA|_V)+\tfrac18\int\mathrm{dvol}\,(d^*A|_V,d^*A|_V)$ \\
total & & $\tfrac94\int\mathrm{dvol}\,(A|_V,A|_V)$ & $-\int\mathrm{dvol}\,(dA|_V,dA|_V)$ \\
\end{tabular}
\tabcaption{Contributions $\hbar|\operatorname{Aut}\gamma|^{-1}\mathrm{Sing}_\varepsilon w_\gamma(\varepsilon)$ of the graphs with external $A|_V,A|_V$ to $I^{CT}[\varepsilon]$, in units of $\hbar\,\operatorname{tr}_{\mathfrak g}(\mathrm{ad}_X\mathrm{ad}_Y)/(16\pi^2(X,Y))$; the $\varepsilon^{-1}$ part multiplies $\varepsilon^{-1}$ and the $\log\varepsilon$ part multiplies $\log\varepsilon$. The coefficients of $\int\mathrm{dvol}\,(d^*A|_V,d^*A|_V)$ cancel between (c) and (f).}

\fig{diagrams/composite/twopoint_AWperp.png}{The one-loop one-particle-irreducible graphs with external $A|_{W_{\mathbb R}^\perp},A|_{W_{\mathbb R}^\perp}$: (a), (b) quartic vertex with an $A|_V$, resp. $A|_{W_{\mathbb R}^\perp}$, loop; (c) bubble with one $A|_V$ and one $A|_{W_{\mathbb R}^\perp}$ line; (d) gaugino loop with two vertices $\lambda\lambda A|_{W_{\mathbb R}^\perp}$. (e): the gaugino loop with external $A|_V,A|_{W_{\mathbb R}^\perp}$.}

\begin{tabular}{l|c|c|c}
graph & $1/|\operatorname{Aut}\gamma|$ & $\varepsilon^{-1}$ part & $\log\varepsilon$ part \\
\hline
(a) & $\tfrac12$ & $2\int\mathrm{dvol}\,(A|_{W_{\mathbb R}^\perp},A|_{W_{\mathbb R}^\perp})$ & $0$ \\
(b) & $\tfrac12$ & $\tfrac52\int\mathrm{dvol}\,(A|_{W_{\mathbb R}^\perp},A|_{W_{\mathbb R}^\perp})$ & $0$ \\
(c) & $1$ & $-\tfrac14\int\mathrm{dvol}\,(A|_{W_{\mathbb R}^\perp},A|_{W_{\mathbb R}^\perp})$ & $\int\mathrm{dvol}\,(dA|_{W_{\mathbb R}^\perp},dA|_{W_{\mathbb R}^\perp})$ \\
(d) & $\tfrac12$ & $-2\int\mathrm{dvol}\,(A|_{W_{\mathbb R}^\perp},A|_{W_{\mathbb R}^\perp})$ & $-2\int\mathrm{dvol}\,(dA|_{W_{\mathbb R}^\perp},dA|_{W_{\mathbb R}^\perp})$ \\
total & & $\tfrac94\int\mathrm{dvol}\,(A|_{W_{\mathbb R}^\perp},A|_{W_{\mathbb R}^\perp})$ & $-\int\mathrm{dvol}\,(dA|_{W_{\mathbb R}^\perp},dA|_{W_{\mathbb R}^\perp})$ \\
(e) & $\tfrac12$ & $0$ & $0$ \\
\end{tabular}
\tabcaption{Contributions of the graphs with external $A|_{W_{\mathbb R}^\perp},A|_{W_{\mathbb R}^\perp}$ (rows (a)--(d)) and $A|_V,A|_{W_{\mathbb R}^\perp}$ (row (e)) to $I^{CT}[\varepsilon]$, in the units of the previous table.}

\fig{diagrams/composite/twopoint_lambda_c_Astar.png}{The one-loop one-particle-irreducible graphs with external $\lambda,\lambda$: (a) with an $A|_V$ line, (b) with an $A|_{W_{\mathbb R}^\perp}$ line; and with external $c,A^*|_V$: (c) with one $A|_V$ line and one $c$--$A^*$ line.}

\begin{tabular}{l|c|c|c}
graph & $1/|\operatorname{Aut}\gamma|$ & $\varepsilon^{-1}$ part & $\log\varepsilon$ part \\
\hline
(a) & $1$ & $0$ & $\int\mathrm{dvol}\,(\lambda,{\not d}\lambda)$ \\
(b) & $1$ & $0$ & $3\int\mathrm{dvol}\,(\lambda,{\not d}\lambda)$ \\
total $\lambda,\lambda$ & & $0$ & $4\int\mathrm{dvol}\,(\lambda,{\not d}\lambda)$ \\
(c) & $1$ & $0$ & $-\int\mathrm{dvol}\,(dc,A^*)$ \\
\end{tabular}
\tabcaption{Contributions of the graphs with external $\lambda,\lambda$ and $c,A^*|_V$ to $I^{CT}[\varepsilon]$, in the units of the previous tables.}

Five representative graphs are computed below in five steps: the amplitude, from the one-loop formula (Proposition \ref{prop:c3-oneloop}), the vertex functions of Definition \ref{def:c1-vertex} and the propagator components of Proposition \ref{prop:c3-P}; the numerator algebra; the Schwinger integral, by Lemma \ref{lem:c4-gauss}; the $\varepsilon$-expansion, by Lemma \ref{lem:c4-integrals}; and the local part, by \eqref{eq:c4-pw}. For two legs, the order-$n$ term of \eqref{eq:c3-oneloop} enters through its part bilinear in the legs: for $n=2$,
\[
\frac\hbar4\mathrm{STr}\big((\text{one-loop operator})^2\big)\Big|_{\text{bilinear}}=\frac\hbar2\mathrm{STr}\big((\text{part linear in leg }1)(\text{part linear in leg }2)\big),
\]
and for $n=1$ it is $\frac\hbar2\mathrm{STr}$ of the part of the one-loop operator bilinear in the legs. In the bubbles the two edges carry the momenta $p_1$ and $p_1-p$, with Schwinger times $t_1$ and $t_2$:
\[
\frac{e^{-\varepsilon|p_1|^2}}{|p_1|^2}=\int_\varepsilon^\infty e^{-t_1|p_1|^2}dt_1,\qquad \frac{e^{-\varepsilon|p_1-p|^2}}{|p_1-p|^2}=\int_\varepsilon^\infty e^{-t_2|p_1-p|^2}dt_2 .
\]
The colour trace of the loop is $\operatorname{tr}_{\mathfrak g}(\mathrm{ad}_{X_1}\mathrm{ad}_{X_2})=\frac{\operatorname{tr}_{\mathfrak g}(\mathrm{ad}_X\mathrm{ad}_Y)}{(X,Y)}(X_1,X_2)$ for the $\mathfrak g$-components $X_1,X_2$ of the two legs, and the factor $(X_1,X_2)$ is part of the pairings $(A(p),A(-p))$ and $(\iota_pA(p),\iota_pA(-p))$.

\begin{example}[Gauge bubble, graph (c) of the $A|_V$ two-point function]
\label{ex:c4-gauge}
Amplitude. The two vertices are \eqref{eq:c1-AAA} at the incoming momenta $(p,-p_1,p_1-p)$ and $(-p,p_1,p-p_1)$; each edge carries the $A|_V\otimes A|_V$ component $2(-,-)^{-1}\otimes(\text{inverse metric of }V)$ of $P(\varepsilon,\infty)$, times $\frac{e^{-\varepsilon|p_1|^2}}{|p_1|^2}$, resp. $\frac{e^{-\varepsilon|p_1-p|^2}}{|p_1-p|^2}$. The plane-wave value of $\hbar|\operatorname{Aut}\gamma|^{-1}w_\gamma(\varepsilon)$ is
\[
\frac\hbar2\,\frac{\operatorname{tr}_{\mathfrak g}(\mathrm{ad}_X\mathrm{ad}_Y)}{(X,Y)}\int\frac{d^4p_1}{(2\pi)^4}\,\frac{e^{-\varepsilon|p_1|^2}}{|p_1|^2}\,\frac{e^{-\varepsilon|p_1-p|^2}}{|p_1-p|^2}\,(\text{numerator}).
\]
Numerator algebra. Contracting the two vertex functions along the two edges (the exchange of the edges is the automorphism counted by $1/|\operatorname{Aut}\gamma|=\tfrac12$),
\[
\begin{aligned}
\text{numerator}={}&-\big(|p_1|^2+|p_1-p|^2+4|p|^2\big)(A(p),A(-p))+2(\iota_pA(p),\iota_pA(-p))\\
&-5\big[(\iota_{p_1}A(p),\iota_{p_1-p}A(-p))+(\iota_{p_1-p}A(p),\iota_{p_1}A(-p))\big].
\end{aligned}
\]
At $p_1=\frac{t_2}{t_1+t_2}p$, $|p_1|^2+|p_1-p|^2=\big(1-\frac{2t_1t_2}{(t_1+t_2)^2}\big)|p|^2$, and each of the two bracketed pairings equals $-\frac{t_1t_2}{(t_1+t_2)^2}(\iota_pA(p),\iota_pA(-p))$; the positive Laplacian in $p_1$ of the numerator is $36\,(A(p),A(-p))$.

Schwinger integral. By Lemma \ref{lem:c4-gauss} the plane-wave value is
\[
\begin{aligned}
\frac\hbar2\,\frac{\operatorname{tr}_{\mathfrak g}(\mathrm{ad}_X\mathrm{ad}_Y)}{16\pi^2(X,Y)}\iint_{[\varepsilon,\infty)^2}\frac{e^{-\frac{t_1t_2|p|^2}{t_1+t_2}}}{(t_1+t_2)^2}\Big[&-\frac{9\,(A(p),A(-p))}{t_1+t_2}+\Big(\frac{2t_1t_2}{(t_1+t_2)^2}-5\Big)|p|^2(A(p),A(-p))\\
&+\Big(\frac{10t_1t_2}{(t_1+t_2)^2}+2\Big)(\iota_pA(p),\iota_pA(-p))\Big]dt_1dt_2 .
\end{aligned}
\]
$\varepsilon$-expansion. By Lemma \ref{lem:c4-integrals} the double integral of the bracket is, modulo $\mathscr A_{\ge0}$,
\[
\begin{aligned}
&-9\Big(\tfrac14\varepsilon^{-1}+\tfrac16|p|^2\log\varepsilon\Big)(A(p),A(-p))+\Big(-\tfrac13+5\Big)|p|^2\log\varepsilon\,(A(p),A(-p))\\
&\quad+\Big(-\tfrac{10}6-2\Big)\log\varepsilon\,(\iota_pA(p),\iota_pA(-p))\\
&\quad=-\tfrac94\varepsilon^{-1}(A(p),A(-p))+\tfrac{19}{6}|p|^2\log\varepsilon\,(A(p),A(-p))-\tfrac{11}{3}\log\varepsilon\,(\iota_pA(p),\iota_pA(-p)).
\end{aligned}
\]
Local part. The plane-wave coefficients are $\tfrac12\cdot(-\tfrac94)=-\tfrac98$, $\tfrac12\cdot\tfrac{19}6=\tfrac{19}{12}$, $\tfrac12\cdot(-\tfrac{11}3)=-\tfrac{11}6$. By \eqref{eq:c4-pw} these are the plane-wave coefficients of
\[
\begin{aligned}
\frac{\hbar\,\operatorname{tr}_{\mathfrak g}(\mathrm{ad}_X\mathrm{ad}_Y)}{16\pi^2(X,Y)}\Big[&-\tfrac9{16}\varepsilon^{-1}\int\mathrm{dvol}\,(A|_V,A|_V)\\
&+\log\varepsilon\Big(\tfrac{19}{24}\int\mathrm{dvol}\,(dA|_V,dA|_V)-\tfrac18\int\mathrm{dvol}\,(d^*A|_V,d^*A|_V)\Big)\Big],
\end{aligned}
\]
since $2\cdot(-\tfrac9{16})=-\tfrac98$, $2\cdot\tfrac{19}{24}=\tfrac{19}{12}$ and $2\big(-\tfrac{19}{24}-\tfrac18\big)=-\tfrac{11}{6}$. This is row (c) of the first table.
\end{example}

\begin{example}[Gaugino loop, graph (e) of the $A|_V$ two-point function]
\label{ex:c4-gaugino}
Amplitude. Two vertices $\lambda\lambda A|_V$ (Definition \ref{def:c1-vertex}); each edge carries the $\lambda\otimes\lambda$ component $(-,-)^{-1}\otimes i\rho(\cdot)|_{\Sigma^*\to\Sigma}$ of $P(\varepsilon,\infty)$ at the momentum $p_1$, resp. $p_1-p$; the parity operator in $\mathrm{STr}$ gives the sign of the loop of odd fields. The prefactor and the scalar propagator factors are those of Example \ref{ex:c4-gauge}.

Numerator algebra. The trace over $\Sigma$ of the products of Clifford multiplications, evaluated with $\rho(v)\rho(v')+\rho(v')\rho(v)=2(v,v')$, gives
\[
\text{numerator}=-16\big[(p_1,p_1-p)(A(p),A(-p))-(\iota_{p_1}A(p),\iota_{p_1-p}A(-p))-(\iota_{p_1-p}A(p),\iota_{p_1}A(-p))\big],
\]
with value $\frac{16t_1t_2}{(t_1+t_2)^2}\big(|p|^2(A(p),A(-p))-2(\iota_pA(p),\iota_pA(-p))\big)$ at the centre and positive Laplacian $64\,(A(p),A(-p))$.

Schwinger integral.
\[
\begin{aligned}
\frac\hbar2\,\frac{\operatorname{tr}_{\mathfrak g}(\mathrm{ad}_X\mathrm{ad}_Y)}{16\pi^2(X,Y)}\iint_{[\varepsilon,\infty)^2}\frac{e^{-\frac{t_1t_2|p|^2}{t_1+t_2}}}{(t_1+t_2)^2}\Big[&-\frac{16\,(A(p),A(-p))}{t_1+t_2}+\frac{16t_1t_2}{(t_1+t_2)^2}|p|^2(A(p),A(-p))\\
&-\frac{32t_1t_2}{(t_1+t_2)^2}(\iota_pA(p),\iota_pA(-p))\Big]dt_1dt_2 .
\end{aligned}
\]
$\varepsilon$-expansion. Modulo $\mathscr A_{\ge0}$, the double integral of the bracket is $-4\varepsilon^{-1}(A(p),A(-p))-\tfrac{16}{3}\log\varepsilon\big(|p|^2(A(p),A(-p))-(\iota_pA(p),\iota_pA(-p))\big)$, the $|p|^2$ terms being $-16\cdot\tfrac16-16\cdot\tfrac16$.

Local part. $\frac{\hbar\,\operatorname{tr}_{\mathfrak g}(\mathrm{ad}_X\mathrm{ad}_Y)}{16\pi^2(X,Y)}\big[-\varepsilon^{-1}\int\mathrm{dvol}\,(A|_V,A|_V)-\tfrac43\log\varepsilon\int\mathrm{dvol}\,(dA|_V,dA|_V)\big]$: the $\log\varepsilon$ part of the gaugino loop is transverse.
\end{example}

\begin{example}[Ghost loop, graph (f) of the $A|_V$ two-point function]
\label{ex:c4-ghost}
Amplitude. Two vertices $A|_V\,c\,A^*|_V$ (Definition \ref{def:c1-vertex}); the edges carry the components $c\otimes A^*|_V$ and $A^*|_V\otimes c$ of $P(\varepsilon,\infty)$, equal to $-(-,-)^{-1}\,i\,(\text{edge momentum})$ times the scalar factor (Proposition \ref{prop:c3-P}); the parity operator gives the sign of the loop of odd fields.

Numerator algebra. The two orientations of the ghost line contribute
\[
\text{numerator}=(\iota_{p_1}A(p),\iota_{p_1-p}A(-p))+(\iota_{p_1-p}A(p),\iota_{p_1}A(-p)),
\]
with value $-\frac{2t_1t_2}{(t_1+t_2)^2}(\iota_pA(p),\iota_pA(-p))$ at the centre and positive Laplacian $-4\,(A(p),A(-p))$.

Schwinger integral.
\[
\frac\hbar2\,\frac{\operatorname{tr}_{\mathfrak g}(\mathrm{ad}_X\mathrm{ad}_Y)}{16\pi^2(X,Y)}\iint_{[\varepsilon,\infty)^2}\frac{e^{-\frac{t_1t_2|p|^2}{t_1+t_2}}}{(t_1+t_2)^2}\Big[\frac{(A(p),A(-p))}{t_1+t_2}-\frac{2t_1t_2}{(t_1+t_2)^2}(\iota_pA(p),\iota_pA(-p))\Big]dt_1dt_2 .
\]
$\varepsilon$-expansion. Modulo $\mathscr A_{\ge0}$, the double integral of the bracket is $\tfrac14\varepsilon^{-1}(A(p),A(-p))+\log\varepsilon\big(\tfrac16|p|^2(A(p),A(-p))+\tfrac13(\iota_pA(p),\iota_pA(-p))\big)$.

Local part.
\[
\begin{aligned}
\frac{\hbar\,\operatorname{tr}_{\mathfrak g}(\mathrm{ad}_X\mathrm{ad}_Y)}{16\pi^2(X,Y)}\Big[&\tfrac1{16}\varepsilon^{-1}\int\mathrm{dvol}\,(A|_V,A|_V)\\
&+\log\varepsilon\Big(\tfrac1{24}\int\mathrm{dvol}\,(dA|_V,dA|_V)+\tfrac18\int\mathrm{dvol}\,(d^*A|_V,d^*A|_V)\Big)\Big],
\end{aligned}
\]
since $2\cdot\tfrac1{16}=\tfrac12\cdot\tfrac14$, $2\cdot\tfrac1{24}=\tfrac12\cdot\tfrac16$ and $2\big(-\tfrac1{24}+\tfrac18\big)=\tfrac12\cdot\tfrac13$. The longitudinal coefficient $\tfrac18$ cancels the $-\tfrac18$ of Example \ref{ex:c4-gauge}.
\end{example}

\begin{example}[Seagull, graph (a) of the $A|_V$ two-point function]
\label{ex:c4-seagull}
Amplitude. One quartic vertex \eqref{eq:c1-AAAA} with two legs and one $A|_V$ edge closed into a loop; the term $n=1$ of \eqref{eq:c3-oneloop}.

Numerator algebra. Contracting two legs of \eqref{eq:c1-AAAA} with $2(-,-)^{-1}\otimes(\text{inverse metric of }V)$ and summing over the pairs of legs gives the numerator $6\,(A(p),A(-p))$, independent of $p_1$.

Schwinger integral. By Lemma \ref{lem:c4-gauss} with one edge,
\[
\frac\hbar2\,\frac{\operatorname{tr}_{\mathfrak g}(\mathrm{ad}_X\mathrm{ad}_Y)}{16\pi^2(X,Y)}\int_\varepsilon^\infty\frac{6\,(A(p),A(-p))}{t^2}\,dt .
\]
$\varepsilon$-expansion. The integral equals $6\,\varepsilon^{-1}(A(p),A(-p))$ exactly; its $\varepsilon\frac{d}{d\varepsilon}$ is the graph with the edge pinned at $t=\varepsilon$, $-\varepsilon\cdot6\,\varepsilon^{-2}(A(p),A(-p))$, as in Proposition \ref{prop:c4-boundary}.

Local part. $\frac{\hbar\,\operatorname{tr}_{\mathfrak g}(\mathrm{ad}_X\mathrm{ad}_Y)}{16\pi^2(X,Y)}\,\tfrac32\,\varepsilon^{-1}\int\mathrm{dvol}\,(A|_V,A|_V)$, a mass term with vanishing $\log\varepsilon$ part.
\end{example}

\begin{example}[The graph with external $c,A^*|_V$]
\label{ex:c4-cAstar}
Amplitude. Legs $r_1A^*(p)e^{i(p,x)}$ and $r_2c(-p)e^{-i(p,x)}$; two vertices $A|_V\,c\,A^*|_V$; one $A|_V$ edge and one $c$--$A^*$ edge (graph (c) of the third figure).

Numerator algebra. The coefficient of $r_1r_2$ is $2i\,(p\,c(-p),A^*(p))$, independent of $p_1$.

Schwinger integral.
\[
\frac\hbar2\,\frac{\operatorname{tr}_{\mathfrak g}(\mathrm{ad}_X\mathrm{ad}_Y)}{16\pi^2(X,Y)}\iint_{[\varepsilon,\infty)^2}\frac{e^{-\frac{t_1t_2|p|^2}{t_1+t_2}}}{(t_1+t_2)^2}\,2i\,(p\,c(-p),A^*(p))\,dt_1dt_2 .
\]
$\varepsilon$-expansion. Modulo $\mathscr A_{\ge0}$, $-2i\log\varepsilon\,(p\,c(-p),A^*(p))$.

Local part. The plane-wave coefficient is $\tfrac12\cdot(-2i)=-i$, and the plane-wave value of $\int\mathrm{dvol}\,(dc,A^*)$ is $i\,(p\,c(-p),A^*(p))$ by \eqref{eq:c4-pw}; the contribution is $\frac{\hbar\,\operatorname{tr}_{\mathfrak g}(\mathrm{ad}_X\mathrm{ad}_Y)}{16\pi^2(X,Y)}\big(-\log\varepsilon\big)\int\mathrm{dvol}\,(dc,A^*)$.
\end{example}

\begin{theorem}[The quadratic counterterm]
\label{thm:c4-quadratic}
The part of $I^{CT}[\varepsilon]$ quadratic in the fields is
\begin{equation}
\begin{aligned}
\frac{\hbar\,\operatorname{tr}_{\mathfrak g}(\mathrm{ad}_X\mathrm{ad}_Y)}{16\pi^2(X,Y)}\Big\{&\tfrac94\,\varepsilon^{-1}\Big[\int\mathrm{dvol}\,(A|_V,A|_V)+\int\mathrm{dvol}\,(A|_{W_{\mathbb R}^\perp},A|_{W_{\mathbb R}^\perp})\Big]\\
&+\log\varepsilon\Big[-\int\mathrm{dvol}\,(dA|_V,dA|_V)-\int\mathrm{dvol}\,(dA|_{W_{\mathbb R}^\perp},dA|_{W_{\mathbb R}^\perp})\\
&\qquad\qquad+4\int\mathrm{dvol}\,(\lambda,{\not d}\lambda)-\int\mathrm{dvol}\,(dc,A^*)\Big]\Big\}.
\end{aligned}
\label{eq:c4-quadratic}
\end{equation}
The coefficient of $\int\mathrm{dvol}\,(d^*A|_V,d^*A|_V)$ is $0$, and every quadratic local functional mixing $A|_V$ with $A|_{W_{\mathbb R}^\perp}$, or $c$ with $A^*|_{W_{\mathbb R}^\perp}$, has coefficient $0$. In terms of the ten-component field $A$,
\[
\begin{aligned}
\frac{\hbar\,\operatorname{tr}_{\mathfrak g}(\mathrm{ad}_X\mathrm{ad}_Y)}{16\pi^2(X,Y)}\Big\{&\tfrac94\,\varepsilon^{-1}\int\mathrm{dvol}\,(A,A)\\
&+\log\varepsilon\Big[-\int\mathrm{dvol}\,(dA,dA)+4\int\mathrm{dvol}\,(\lambda,{\not d}\lambda)-\int\mathrm{dvol}\,(dc,A^*)\Big]\Big\}.
\end{aligned}
\]
\end{theorem}

\begin{proof}
$I^{CT}[\varepsilon]$ is even and of degree $0$, so by the degrees of Chapter 1 its quadratic part has legs $(A,A)$, $(\lambda,\lambda)$, $(c,A^*)$ or $(c,\lambda^*)$. The regulator is covariant under $\mathrm{Spin}(4,\mathbb C)\times\mathrm{Spin}(6,\mathbb C)$, since $H=d^*d+dd^*$ and $Q^{GF}$ commute with it, so $I^{CT}[\varepsilon]$ is invariant. The invariant pairings of a vector of $V$ with a vector of $W_{\mathbb R}^\perp$, of a scalar with a spinor, and of a scalar with a vector of $W_{\mathbb R}^\perp$ vanish; this gives the vanishing mixings and the vanishing of the sectors $(c,\lambda^*)$ and $(c,A^*|_{W_{\mathbb R}^\perp})$. The remaining sectors are the sums of the three tables, whose rows are all one-particle-irreducible one-loop graphs with these legs; the representative rows are derived in Examples \ref{ex:c4-gauge}--\ref{ex:c4-cAstar}. The sums are
\[
\begin{aligned}
A|_V:&\quad \tfrac32+3-\tfrac9{16}-\tfrac34-1+\tfrac1{16}=\tfrac94,\qquad \tfrac{19}{24}-\tfrac12-\tfrac43+\tfrac1{24}=-1,\qquad -\tfrac18+\tfrac18=0;\\
A|_{W_{\mathbb R}^\perp}:&\quad 2+\tfrac52-\tfrac14-2=\tfrac94,\qquad 1-2=-1;\qquad \lambda:\quad 1+3=4;\qquad (c,A^*):\quad -1 .
\end{aligned}
\]
The ten-component form uses $(dA,dA)=(dA|_V,dA|_V)+(dA|_{W_{\mathbb R}^\perp},dA|_{W_{\mathbb R}^\perp})$, the cross term vanishing by the orthogonality of $V$ and $W_{\mathbb R}^\perp$.
\end{proof}

\begin{remark}[Two properties of the quadratic counterterm]
\label{rem:c4-quadratic}
The per-graph entries of the $A|_V$ and $A|_{W_{\mathbb R}^\perp}$ sectors differ, and their totals agree: \eqref{eq:c4-quadratic} is a functional of the ten-component field $A$. The $\log\varepsilon$ part is transverse, the longitudinal coefficients $-\tfrac18$ of the gauge bubble and $+\tfrac18$ of the ghost loop cancelling, and it is a sum of multiples of the three terms of the free action \eqref{eq:c1-free}:
\[
\begin{aligned}
&-\int\mathrm{dvol}\,(dA,dA)+4\int\mathrm{dvol}\,(\lambda,{\not d}\lambda)-\int\mathrm{dvol}\,(dc,A^*)\\
&\quad=4\Big(-\tfrac14\int\mathrm{dvol}\,(dA,dA)\Big)+8\Big(\tfrac12\int\mathrm{dvol}\,(\lambda,{\not d}\lambda)\Big)-1\cdot\int\mathrm{dvol}\,(dc,A^*).
\end{aligned}
\]
The $\varepsilon^{-1}$ part is a mass term for all ten components of $A$.
\end{remark}

\section{Cubic and quartic graphs}

The cubic and quartic parts of $I^{CT}[\varepsilon]$ are organised by the Weyl weights of Example \ref{ex:c2-dil}: $A$ has weight $1$, $\lambda$ weight $\tfrac32$, $c$ weight $0$, and the antifields $A^*,\lambda^*,c^*$ the complementary weights $3,\tfrac52,4$, so that each pairing of a field with its antifield has weight $4$. The weight of a local functional $\int\mathrm{dvol}\,(\cdots)$ is the sum of the weights of its fields plus the number of its derivatives, and every term of $S$ has weight $4$. The superficial degree of a graph is $4$ minus the sum of the weights of its legs: $2$ for $(A,A)$, $1$ for $(\lambda,\lambda)$, $(c,A^*)$, $(A|_V,A|_V,A|_V)$, $(A|_V,A|_{W_{\mathbb R}^\perp},A|_{W_{\mathbb R}^\perp})$, and $0$ for the other cubic and the quartic leg sets below.

\begin{lemma}[Weights of the singular coefficients]
\label{lem:c4-weights}
Let $\gamma$ be a one-loop one-particle-irreducible graph with vertices labelled by $I$. Its superficial degree in the sense of Lemma \ref{lem:c4-local} is $4$ minus the sum of the weights of its legs. The coefficient of $\varepsilon^{-k}$ in its singular part is a local functional of weight $4-2k$, and the coefficient of $\log\varepsilon$ a local functional of weight $4$; in particular the coefficient of $\varepsilon^{-1}$ has weight $2$.
\end{lemma}

\begin{proof}
In momentum space, a vertex with fields of total weight $4-k$ is a polynomial of degree $k$ in the momenta, since every term of $S$ has weight $4$; an edge component of $P(\varepsilon,\infty)$ between slots of total weight $2+k$ is a polynomial of degree $k$ times $\int_\varepsilon^\infty e^{-t_e|\text{momentum}|^2}dt_e$ (Proposition \ref{prop:c3-P}: degree $0$ for $A\otimes A$, $1$ for $c\otimes A^*$ and $\lambda\otimes\lambda$). Each internal half-edge contributes its weight once to a vertex and once to an edge, so, for legs in fixed slots and every assignment of slots to the internal half-edges, the numerator of $\gamma$ is homogeneous of degree $4|V(\gamma)|-2|E(\gamma)|$ minus the sum of the weights of the legs, in the loop momentum $p_1$ and the leg momenta, which is the homogeneity required in Lemma \ref{lem:c4-local}; with $|E(\gamma)|=|V(\gamma)|$ at one loop, its degree plus $4-2|E(\gamma)|$ is $4$ minus the sum of the weights of the legs. By Lemma \ref{lem:c4-local} the coefficient of $\varepsilon^{-k}$ is a polynomial in the leg momenta, homogeneous of degree (superficial degree)$-2k$, and that of $\log\varepsilon$ is homogeneous of degree equal to the superficial degree. The weight of such a local coefficient is the sum of the leg weights, $4-$(superficial degree), plus its degree in the momenta, which gives $4-2k$ and $4$.
\end{proof}

\begin{lemma}[Local functionals of weight at most four]
\label{lem:c4-enumeration}
Let a local functional of degree $0$, without constant term, have weight at most $4$ and be invariant under $\mathrm{Spin}(4,\mathbb C)\times\mathrm{Spin}(6,\mathbb C)$ and the adjoint action of $\mathfrak g$. If it is even, it is a linear combination of: in weight $2$, $\int\mathrm{dvol}\,(A|_V,A|_V)$ and $\int\mathrm{dvol}\,(A|_{W_{\mathbb R}^\perp},A|_{W_{\mathbb R}^\perp})$; in weight $4$, functionals whose monomials have one of the field contents
\[
\begin{aligned}
&(A|_V)^2\ \text{and}\ (A|_{W_{\mathbb R}^\perp})^2\ \text{with two derivatives};\qquad \lambda\lambda\ \text{and}\ c\,A^*|_V\ \text{with one derivative};\\
&(A|_V)^3\ \text{and}\ A|_V(A|_{W_{\mathbb R}^\perp})^2\ \text{with one derivative};\qquad \lambda\lambda A|_V,\ \lambda\lambda A|_{W_{\mathbb R}^\perp},\\
&c\,A|_V\,A^*|_V,\ c\,A|_{W_{\mathbb R}^\perp}A^*|_{W_{\mathbb R}^\perp},\ c\,\lambda\,\lambda^*,\ c\,c\,c^*;\qquad (A|_V)^4,\ (A|_V)^2(A|_{W_{\mathbb R}^\perp})^2,\ (A|_{W_{\mathbb R}^\perp})^4,
\end{aligned}
\]
In weights $1$ and $3$ there are none.
\end{lemma}

\begin{proof}
As functionals, $c$ has degree $+1$, $A^*$ and $\lambda^*$ degree $-1$, $c^*$ degree $-2$, and $A,\lambda$ degree $0$ (Chapter 1). Degree $0$ requires the number of $c$ to equal the number of $A^*$ and $\lambda^*$ plus twice the number of $c^*$, and weight at most $4$ allows at most one antifield. With $c^*$ (weight $4$): $c\,c\,c^*$. With $A^*$ (weight $3$): one $c$ and weight at most $1$ more, i.e. $c\,A^*$ with one derivative or $c\,A\,A^*$; invariance pairs $A^*|_V$ with the derivative or with $A|_V$, and $A^*|_{W_{\mathbb R}^\perp}$ with $A|_{W_{\mathbb R}^\perp}$. With $\lambda^*$ (weight $\tfrac52$): one $c$ and, by $\mathrm{Spin}(4,\mathbb C)$, one spinor $\lambda$. Without antifields there is no $c$; the monomials are built from $A$, $\lambda$ and derivatives, with an even number of spinors. $\mathrm{Spin}(6,\mathbb C)$ admits invariants of an even number of vectors of $W_{\mathbb R}^\perp$ (the invariant volume form of $W_{\mathbb R}^\perp$ needs six) and $\mathrm{Spin}(4,\mathbb C)$ of an even number of vector indices from $A|_V$ and the derivatives. In weight $2$ this leaves $(A|_V)^2$, $(A|_{W_{\mathbb R}^\perp})^2$; in weight $4$ the listed contents, a single $A$ with three derivatives being a total derivative; in weight $1$ the single $A$, without invariant; in weight $3$ the contents $A^3$, $A\,A$ with one derivative, $c\,A^*$, each with an odd number of vector indices of $V$ or of $W_{\mathbb R}^\perp$, and $\lambda\lambda$, whose components lie in $S_+\otimes S_-$ (a vector of $V$) or in $W\otimes W$, $W^*\otimes W^*$ (without $SL(4,\mathbb C)$-invariant); all of these have vanishing invariants.
\end{proof}

\begin{proposition}[Structure of the one-loop counterterm]
\label{prop:c4-structure}
$I^{CT}[\varepsilon]$ is the sum of $\varepsilon^{-1}$ times the mass term of Theorem \ref{thm:c4-quadratic} and $\log\varepsilon$ times a local functional of weight $4$ with the field contents of Lemma \ref{lem:c4-enumeration}. Its $\log\varepsilon$ part is therefore determined by the one-particle-irreducible one-loop graphs with the leg sets $(A|_V,A|_V)$, $(A|_{W_{\mathbb R}^\perp},A|_{W_{\mathbb R}^\perp})$, $(\lambda,\lambda)$, $(c,A^*|_V)$, the eight cubic sets $(A|_V,A|_V,A|_V)$, $(A|_V,A|_{W_{\mathbb R}^\perp},A|_{W_{\mathbb R}^\perp})$, $(\lambda,\lambda,A|_V)$, $(\lambda,\lambda,A|_{W_{\mathbb R}^\perp})$, $(c,A|_V,A^*|_V)$, $(c,A|_{W_{\mathbb R}^\perp},A^*|_{W_{\mathbb R}^\perp})$, $(c,\lambda,\lambda^*)$, $(c,c,c^*)$, and the three quartic sets $(A|_V)^4$, $(A|_V)^2(A|_{W_{\mathbb R}^\perp})^2$, $(A|_{W_{\mathbb R}^\perp})^4$.
\end{proposition}

\begin{proof}
By Lemmas \ref{lem:c4-local} and \ref{lem:c4-weights}, the singular part of a one-particle-irreducible one-loop graph with vertices labelled by $I$ consists of $\varepsilon^{-k}$ times a local functional of weight $4-2k$, $1\le k\le\frac12(\text{superficial degree})$, and $\log\varepsilon$ times a local functional of weight $4$, and it vanishes when the superficial degree is negative. The slots of every vertex of $I$ have degrees adding to $0$, and so have the two slots joined by every nonzero component of $P(\varepsilon,\infty)$ (Proposition \ref{prop:c3-P}: $A\otimes A$, $c\otimes A^*|_V$, $A^*|_V\otimes c$, $\lambda\otimes\lambda$); hence the legs of every graph have degrees adding to $0$. With the degrees of the proof of Lemma \ref{lem:c4-enumeration}, a set of $n$ legs with degrees adding to $0$ has weight at least $n$: every leg other than $c$ has weight at least $1$, every $c$ has weight $0$ and degree $+1$, and the antifields exceed weight $1$ by $\tfrac32$ per unit of negative degree or more ($A^*$: $3=1+2$, $\lambda^*$: $\tfrac52=1+\tfrac32$, $c^*$: $4=1+\tfrac32\cdot2$), so the weight is at least $n-(\text{number of }c)+\tfrac32(\text{number of }c)\ge n$. The graphs with one leg have value $0$ (Lemma \ref{lem:c4-factor}). Hence the superficial degree, $4$ minus the weight of the legs, is at most $2$, and negative for five or more legs; by Lemma \ref{lem:c4-local} no power $\varepsilon^{-k}$ with $k\ge2$ occurs, and the coefficient of $\varepsilon^{-1}$ has weight $2$. $I^{CT}[\varepsilon]$ is invariant under $\mathrm{Spin}(4,\mathbb C)\times\mathrm{Spin}(6,\mathbb C)$ and the adjoint action of $\mathfrak g$ (proof of Theorem \ref{thm:c4-quadratic}), so by Lemma \ref{lem:c4-enumeration} the coefficient of $\varepsilon^{-1}$ is a combination of the two mass terms, computed in Theorem \ref{thm:c4-quadratic}, and the coefficient of $\log\varepsilon$ has the listed contents.
\end{proof}

\subsection{Cubic graphs}

Every cubic graph below has vanishing coefficient of $\varepsilon^{-1}$, in accordance with Lemma \ref{lem:c4-weights}. Its $\log\varepsilon$ coefficient, as a multilinear form in the three legs, is a rational multiple of the tree vertex function of Definition \ref{def:c1-vertex} with the same legs; the tables give these multiples, in units of $\hbar\,\operatorname{tr}_{\mathfrak g}(\mathrm{ad}_X\mathrm{ad}_Y)/(16\pi^2(X,Y))$. The tree vertices are: the $(A|_V)^3$ and $A|_V(A|_{W_{\mathbb R}^\perp})^2$ parts of $\int\mathrm{dvol}\,\big(-\tfrac12(dA,\tfrac12[A\wedge A])\big)$; the $A|_V$ and $A|_{W_{\mathbb R}^\perp}$ parts of $\int\mathrm{dvol}\,\tfrac12(\lambda,({\not d}_A-{\not d})\lambda)$; the $A|_V$ and $A|_{W_{\mathbb R}^\perp}$ parts of $\int\mathrm{dvol}\,([A,c],A^*)$; $\int\mathrm{dvol}\,([c,\lambda],\lambda^*)$; $\int\mathrm{dvol}\,\tfrac12([c,c],c^*)$.

\fig{diagrams/composite/cubic_AV_AV_AV.png}{One-loop one-particle-irreducible graphs with external $(A|_V,A|_V,A|_V)$: triangles with (a) $A|_V$, (b) $A|_{W_{\mathbb R}^\perp}$, (c) $\lambda$ lines, (d), (e) the two orientations of the $c$--$A^*$ loop; bubbles with one quartic vertex and (f)--(h) $A|_V$, (i)--(k) $A|_{W_{\mathbb R}^\perp}$ lines, one for each leg at the cubic vertex.}

\begin{tabular}{l|l|l|c|c|c}
graphs & vertices & lines & number & $1/|\operatorname{Aut}\gamma|$ & $\log\varepsilon$ coefficient \\
\hline
(a) & $(A|_V)^3$, $(A|_V)^3$, $(A|_V)^3$ & $A|_V$, $A|_V$, $A|_V$ & $1$ & $1$ & $\tfrac{13}{4}$ \\
(b) & $A|_V(A|_{W_{\mathbb R}^\perp})^2$, $A|_V(A|_{W_{\mathbb R}^\perp})^2$, $A|_V(A|_{W_{\mathbb R}^\perp})^2$ & $A|_{W_{\mathbb R}^\perp}$, $A|_{W_{\mathbb R}^\perp}$, $A|_{W_{\mathbb R}^\perp}$ & $1$ & $1$ & $2$ \\
(c) & $\lambda\lambda A|_V$, $\lambda\lambda A|_V$, $\lambda\lambda A|_V$ & $\lambda$, $\lambda$, $\lambda$ & $1$ & $1$ & $\tfrac{16}{3}$ \\
(d), (e) & $A|_V\,c\,A^*|_V$, $A|_V\,c\,A^*|_V$, $A|_V\,c\,A^*|_V$ & $c$--$A^*$, $c$--$A^*$, $c$--$A^*$ & $2$ & $1$ & $-\tfrac{1}{12}$ \\
(f)--(h) & $(A|_V)^3$, $(A|_V)^4$ & $A|_V$, $A|_V$ & $3$ & $\tfrac12$ & $-\tfrac{9}{2}$ \\
(i)--(k) & $A|_V(A|_{W_{\mathbb R}^\perp})^2$, $(A|_V)^2(A|_{W_{\mathbb R}^\perp})^2$ & $A|_{W_{\mathbb R}^\perp}$, $A|_{W_{\mathbb R}^\perp}$ & $3$ & $\tfrac12$ & $0$ \\
total & & & & & $6$ \\
\end{tabular}
\tabcaption{The $\log\varepsilon$ coefficients of the graphs with external $(A|_V,A|_V,A|_V)$, as multiples of the tree vertex.}

\fig{diagrams/composite/cubic_AV_AWperp_AWperp.png}{One-loop one-particle-irreducible graphs with external $(A|_V,A|_{W_{\mathbb R}^\perp},A|_{W_{\mathbb R}^\perp})$: triangles (a)--(c); bubbles with one quartic vertex (d)--(g).}

\begin{tabular}{l|l|l|c|c|c}
graphs & vertices & lines & number & $1/|\operatorname{Aut}\gamma|$ & $\log\varepsilon$ coefficient \\
\hline
(a) & $(A|_V)^3$, $A|_V(A|_{W_{\mathbb R}^\perp})^2$, $A|_V(A|_{W_{\mathbb R}^\perp})^2$ & $A|_V$, $A|_V$, $A|_{W_{\mathbb R}^\perp}$ & $1$ & $1$ & $\tfrac{3}{2}$ \\
(b) & $A|_V(A|_{W_{\mathbb R}^\perp})^2$, $A|_V(A|_{W_{\mathbb R}^\perp})^2$, $A|_V(A|_{W_{\mathbb R}^\perp})^2$ & $A|_V$, $A|_{W_{\mathbb R}^\perp}$, $A|_{W_{\mathbb R}^\perp}$ & $1$ & $1$ & $1$ \\
(c) & $\lambda\lambda A|_V$, $\lambda\lambda A|_{W_{\mathbb R}^\perp}$, $\lambda\lambda A|_{W_{\mathbb R}^\perp}$ & $\lambda$, $\lambda$, $\lambda$ & $1$ & $1$ & $8$ \\
(d) & $(A|_V)^3$, $(A|_V)^2(A|_{W_{\mathbb R}^\perp})^2$ & $A|_V$, $A|_V$ & $1$ & $\tfrac12$ & $0$ \\
(e) & $A|_V(A|_{W_{\mathbb R}^\perp})^2$, $(A|_{W_{\mathbb R}^\perp})^4$ & $A|_{W_{\mathbb R}^\perp}$, $A|_{W_{\mathbb R}^\perp}$ & $1$ & $\tfrac12$ & $0$ \\
(f), (g) & $A|_V(A|_{W_{\mathbb R}^\perp})^2$, $(A|_V)^2(A|_{W_{\mathbb R}^\perp})^2$ & $A|_V$, $A|_{W_{\mathbb R}^\perp}$ & $2$ & $1$ & $-\tfrac{9}{2}$ \\
total & & & & & $6$ \\
\end{tabular}
\tabcaption{The $\log\varepsilon$ coefficients of the graphs with external $(A|_V,A|_{W_{\mathbb R}^\perp},A|_{W_{\mathbb R}^\perp})$, as multiples of the tree vertex.}

\fig{diagrams/composite/cubic_lambda_lambda_AV.png}{One-loop one-particle-irreducible graphs with external $(\lambda,\lambda,A|_V)$.}

\begin{tabular}{l|l|l|c|c|c}
graphs & vertices & lines & number & $1/|\operatorname{Aut}\gamma|$ & $\log\varepsilon$ coefficient \\
\hline
(a) & $(A|_V)^3$, $\lambda\lambda A|_V$, $\lambda\lambda A|_V$ & $A|_V$, $A|_V$, $\lambda$ & $1$ & $1$ & $3$ \\
(b) & $A|_V(A|_{W_{\mathbb R}^\perp})^2$, $\lambda\lambda A|_{W_{\mathbb R}^\perp}$, $\lambda\lambda A|_{W_{\mathbb R}^\perp}$ & $A|_{W_{\mathbb R}^\perp}$, $A|_{W_{\mathbb R}^\perp}$, $\lambda$ & $1$ & $1$ & $3$ \\
(c) & $\lambda\lambda A|_V$, $\lambda\lambda A|_V$, $\lambda\lambda A|_V$ & $A|_V$, $\lambda$, $\lambda$ & $1$ & $1$ & $1$ \\
(d) & $\lambda\lambda A|_V$, $\lambda\lambda A|_{W_{\mathbb R}^\perp}$, $\lambda\lambda A|_{W_{\mathbb R}^\perp}$ & $A|_{W_{\mathbb R}^\perp}$, $\lambda$, $\lambda$ & $1$ & $1$ & $3$ \\
total & & & & & $10$ \\
\end{tabular}
\tabcaption{The $\log\varepsilon$ coefficients of the graphs with external $(\lambda,\lambda,A|_V)$, as multiples of the tree vertex; the tree vertex pairs the two four-dimensional chiralities $S_+\otimes W$ and $S_-\otimes W^*$ of $\lambda$.}

\fig{diagrams/composite/cubic_lambda_lambda_AWperp.png}{One-loop one-particle-irreducible graphs with external $(\lambda,\lambda,A|_{W_{\mathbb R}^\perp})$; (a), (b) differ by the leg attached to the vertex $A|_V(A|_{W_{\mathbb R}^\perp})^2$.}

\begin{tabular}{l|l|l|c|c|c}
graphs & vertices & lines & number & $1/|\operatorname{Aut}\gamma|$ & $\log\varepsilon$ coefficient \\
\hline
(a), (b) & $A|_V(A|_{W_{\mathbb R}^\perp})^2$, $\lambda\lambda A|_V$, $\lambda\lambda A|_{W_{\mathbb R}^\perp}$ & $A|_V$, $A|_{W_{\mathbb R}^\perp}$, $\lambda$ & $2$ & $1$ & $2$ \\
(c) & $\lambda\lambda A|_V$, $\lambda\lambda A|_V$, $\lambda\lambda A|_{W_{\mathbb R}^\perp}$ & $A|_V$, $\lambda$, $\lambda$ & $1$ & $1$ & $4$ \\
(d) & $\lambda\lambda A|_{W_{\mathbb R}^\perp}$, $\lambda\lambda A|_{W_{\mathbb R}^\perp}$, $\lambda\lambda A|_{W_{\mathbb R}^\perp}$ & $A|_{W_{\mathbb R}^\perp}$, $\lambda$, $\lambda$ & $1$ & $1$ & $4$ \\
total & & & & & $10$ \\
\end{tabular}
\tabcaption{The $\log\varepsilon$ coefficients of the graphs with external $(\lambda,\lambda,A|_{W_{\mathbb R}^\perp})$, as multiples of the tree vertex. The tree vertex has two structures, both $\lambda$ in $S_+\otimes W$ or both in $S_-\otimes W^*$ (Proposition \ref{prop:c1-chirality}); the total is $10$ times the tree in each of them.}

\fig{diagrams/composite/cubic_c_AV_AstarV.png}{One-loop one-particle-irreducible graphs with external $(c,A|_V,A^*|_V)$.}

\begin{tabular}{l|l|l|c|c|c}
graphs & vertices & lines & number & $1/|\operatorname{Aut}\gamma|$ & $\log\varepsilon$ coefficient \\
\hline
(a) & $(A|_V)^3$, $A|_V\,c\,A^*|_V$, $A|_V\,c\,A^*|_V$ & $A|_V$, $A|_V$, $c$--$A^*$ & $1$ & $1$ & $\tfrac{3}{4}$ \\
(b) & $A|_V\,c\,A^*|_V$, $A|_V\,c\,A^*|_V$, $A|_V\,c\,A^*|_V$ & $A|_V$, $c$--$A^*$, $c$--$A^*$ & $1$ & $1$ & $\tfrac{1}{4}$ \\
total & & & & & $1$ \\
\end{tabular}
\tabcaption{The $\log\varepsilon$ coefficients of the graphs with external $(c,A|_V,A^*|_V)$, as multiples of the tree vertex.}

\fig{diagrams/composite/cubic_c_AWperp_AstarWperp.png}{The one-loop one-particle-irreducible graph with external $(c,A|_{W_{\mathbb R}^\perp},A^*|_{W_{\mathbb R}^\perp})$.}

\begin{tabular}{l|l|l|c|c|c}
graphs & vertices & lines & number & $1/|\operatorname{Aut}\gamma|$ & $\log\varepsilon$ coefficient \\
\hline
(a) & $A|_V(A|_{W_{\mathbb R}^\perp})^2$, $A|_V\,c\,A^*|_V$, $A|_{W_{\mathbb R}^\perp}\,c\,A^*|_{W_{\mathbb R}^\perp}$ & $A|_V$, $A|_{W_{\mathbb R}^\perp}$, $c$--$A^*$ & $1$ & $1$ & $1$ \\
total & & & & & $1$ \\
\end{tabular}
\tabcaption{The $\log\varepsilon$ coefficient of the graph with external $(c,A|_{W_{\mathbb R}^\perp},A^*|_{W_{\mathbb R}^\perp})$, as a multiple of the tree vertex.}

\fig{diagrams/composite/cubic_c_lambda_lambdastar.png}{The one-loop one-particle-irreducible graph with external $(c,\lambda,\lambda^*)$.}

\begin{tabular}{l|l|l|c|c|c}
graphs & vertices & lines & number & $1/|\operatorname{Aut}\gamma|$ & $\log\varepsilon$ coefficient \\
\hline
(a) & $c\,\lambda\,\lambda^*$, $A|_V\,c\,A^*|_V$, $\lambda\lambda A|_V$ & $A|_V$, $c$--$A^*$, $\lambda$ & $1$ & $1$ & $1$ \\
total & & & & & $1$ \\
\end{tabular}
\tabcaption{The $\log\varepsilon$ coefficient of the graph with external $(c,\lambda,\lambda^*)$, as a multiple of the tree vertex, for $\lambda$ in either chirality.}

\fig{diagrams/composite/cubic_c_c_cstar.png}{The one-loop one-particle-irreducible graph with external $(c,c,c^*)$.}

\begin{tabular}{l|l|l|c|c|c}
graphs & vertices & lines & number & $1/|\operatorname{Aut}\gamma|$ & $\log\varepsilon$ coefficient \\
\hline
(a) & $c\,c\,c^*$, $A|_V\,c\,A^*|_V$, $A|_V\,c\,A^*|_V$ & $A|_V$, $c$--$A^*$, $c$--$A^*$ & $1$ & $1$ & $1$ \\
total & & & & & $1$ \\
\end{tabular}
\tabcaption{The $\log\varepsilon$ coefficient of the graph with external $(c,c,c^*)$, as a multiple of the tree vertex.}

\subsection{Quartic graphs}

For the quartic sets leg $j\in\{1,2,3,4\}$ carries a decomposable fibre element $(v_j,-)\otimes X_j$ with $X_j\in\mathfrak g$, $v_j\in V$ for a leg $A|_V$ and $v_j\in W_{\mathbb R}^\perp$ for a leg $A|_{W_{\mathbb R}^\perp}$; in the set $(A|_V)^2(A|_{W_{\mathbb R}^\perp})^2$ legs $1,2$ are $A|_V$ and legs $3,4$ are $A|_{W_{\mathbb R}^\perp}$. A multilinear form is determined by its values on such legs, and by Lemma \ref{lem:c4-weights} the $\log\varepsilon$ coefficients are independent of the momenta. They are written in the basis of products of a colour factor and a Lorentz factor:
\begin{itemize}
\item colour factors: the pairing $([X_1,X_2],[X_3,X_4])$ of the bracket of the colour elements of legs $1,2$ with the bracket of those of legs $3,4$; the pairing $([X_1,X_4],[X_2,X_3])$; and the averaged trace $\tfrac1{24}\sum_{\text{orderings}}\operatorname{tr}_{\mathfrak g}(\mathrm{ad}_{X_1}\mathrm{ad}_{X_2}\mathrm{ad}_{X_3}\mathrm{ad}_{X_4})$, the fourth-order adjoint trace averaged over the $24$ orderings of $X_1,\dots,X_4$. The third pairing $([X_1,X_3],[X_2,X_4])$ is the sum of the first two, by the Jacobi identity and the invariance of $(-,-)$.
\item Lorentz factors: $(v_1,v_2)(v_3,v_4)$, $(v_1,v_3)(v_2,v_4)$, $(v_1,v_4)(v_2,v_3)$, products of metric pairings.
\end{itemize}
The two bracket pairings carry the unit $\hbar\,\operatorname{tr}_{\mathfrak g}(\mathrm{ad}_X\mathrm{ad}_Y)/(16\pi^2(X,Y))$, the averaged trace the unit $\hbar/16\pi^2$. On such legs, by the Hodge normalisation of Definition \ref{def:c1-action},
\[
\begin{aligned}
&([A(p_1)\wedge A(p_2)],[A(p_3)\wedge A(p_4)])\\
&\quad=([X_1,X_2],[X_3,X_4])\big((v_1,v_3)(v_2,v_4)-(v_1,v_4)(v_2,v_3)\big),
\end{aligned}
\]
and similarly for the other two terms of \eqref{eq:c1-AAAA}. With the Jacobi identity in the form
\[
([X_1,X_3],[X_2,X_4])=([X_1,X_2],[X_3,X_4])+([X_1,X_4],[X_2,X_3]),
\]
the tree vertex function \eqref{eq:c1-AAAA} is
\begin{equation}
\begin{aligned}
-\tfrac12\Big[&([X_1,X_2],[X_3,X_4])\big((v_1,v_2)(v_3,v_4)+(v_1,v_3)(v_2,v_4)-2(v_1,v_4)(v_2,v_3)\big)\\
&+([X_1,X_4],[X_2,X_3])\big(2(v_1,v_2)(v_3,v_4)-(v_1,v_3)(v_2,v_4)-(v_1,v_4)(v_2,v_3)\big)\Big],
\end{aligned}
\label{eq:c4-tree4}
\end{equation}
and for $(A|_V)^2(A|_{W_{\mathbb R}^\perp})^2$, where $(v_1,v_3)=(v_1,v_4)=0$, it is
\[
-\tfrac12\big[([X_1,X_2],[X_3,X_4])+2([X_1,X_4],[X_2,X_3])\big](v_1,v_2)(v_3,v_4).
\]

\fig{diagrams/composite/quartic_AV4_boxes.png}{Box graphs with external $(A|_V)^4$. Columns: the three cyclic orders of the legs. Rows: loops of $A|_V$ (a)--(c), $A|_{W_{\mathbb R}^\perp}$ (d)--(f), $\lambda$ (g)--(i), and $c$--$A^*$ in the two orientations (j)--(o).}

\fig{diagrams/composite/quartic_AV4_triangles.png}{Triangle graphs (one quartic vertex) with external $(A|_V)^4$: loops of $A|_V$ (a)--(f) and of $A|_{W_{\mathbb R}^\perp}$ (g)--(l).}

\fig{diagrams/composite/quartic_AV4_bubbles.png}{Bubble graphs (two quartic vertices) with external $(A|_V)^4$: $A|_V$ lines (a)--(c), $A|_{W_{\mathbb R}^\perp}$ lines (d)--(f).}

\begin{tabular}{l|l|l|c|c}
graphs & vertices & lines & number & $1/|\operatorname{Aut}\gamma|$ \\
\hline
box (a)--(c) & $(A|_V)^3$, $(A|_V)^3$, $(A|_V)^3$, $(A|_V)^3$ & $A|_V$, $A|_V$, $A|_V$, $A|_V$ & $3$ & $1$ \\
box (d)--(f) & $A|_V(A|_{W_{\mathbb R}^\perp})^2$, $A|_V(A|_{W_{\mathbb R}^\perp})^2$, $A|_V(A|_{W_{\mathbb R}^\perp})^2$, $A|_V(A|_{W_{\mathbb R}^\perp})^2$ & $A|_{W_{\mathbb R}^\perp}$, $A|_{W_{\mathbb R}^\perp}$, $A|_{W_{\mathbb R}^\perp}$, $A|_{W_{\mathbb R}^\perp}$ & $3$ & $1$ \\
box (g)--(i) & $\lambda\lambda A|_V$, $\lambda\lambda A|_V$, $\lambda\lambda A|_V$, $\lambda\lambda A|_V$ & $\lambda$, $\lambda$, $\lambda$, $\lambda$ & $3$ & $1$ \\
box (j)--(o) & $A|_V\,c\,A^*|_V$, $A|_V\,c\,A^*|_V$, $A|_V\,c\,A^*|_V$, $A|_V\,c\,A^*|_V$ & $c$--$A^*$, $c$--$A^*$, $c$--$A^*$, $c$--$A^*$ & $6$ & $1$ \\
triangle (a)--(f) & $(A|_V)^3$, $(A|_V)^3$, $(A|_V)^4$ & $A|_V$, $A|_V$, $A|_V$ & $6$ & $1$ \\
triangle (g)--(l) & $A|_V(A|_{W_{\mathbb R}^\perp})^2$, $A|_V(A|_{W_{\mathbb R}^\perp})^2$, $(A|_V)^2(A|_{W_{\mathbb R}^\perp})^2$ & $A|_{W_{\mathbb R}^\perp}$, $A|_{W_{\mathbb R}^\perp}$, $A|_{W_{\mathbb R}^\perp}$ & $6$ & $1$ \\
bubble (a)--(c) & $(A|_V)^4$, $(A|_V)^4$ & $A|_V$, $A|_V$ & $3$ & $\tfrac12$ \\
bubble (d)--(f) & $(A|_V)^2(A|_{W_{\mathbb R}^\perp})^2$, $(A|_V)^2(A|_{W_{\mathbb R}^\perp})^2$ & $A|_{W_{\mathbb R}^\perp}$, $A|_{W_{\mathbb R}^\perp}$ & $3$ & $\tfrac12$ \\
\end{tabular}
\tabcaption{The classes of one-loop graphs with external $(A|_V)^4$; the panel letters refer to the box, triangle and bubble figures of this set.}

\begin{tabular}{l|l|c|c|c|c|c|c|c|c|c|c}
colour & Lorentz & box (a)--(c) & box (d)--(f) & box (g)--(i) & box (j)--(o) & triangle (a)--(f) & triangle (g)--(l) & bubble (a)--(c) & bubble (d)--(f) & total & $8\times$tree \\
\hline
$([X_1,X_2],[X_3,X_4])$ & $(v_1,v_2)(v_3,v_4)$ & $\tfrac{5}{12}$ & $0$ & $-\tfrac{8}{3}$ & $0$ & $-\tfrac{13}{4}$ & $-2$ & $\tfrac{5}{2}$ & $1$ & $-4$ & $-4$ \\
$([X_1,X_2],[X_3,X_4])$ & $(v_1,v_3)(v_2,v_4)$ & $\tfrac{5}{12}$ & $0$ & $-\tfrac{8}{3}$ & $0$ & $-\tfrac{13}{4}$ & $-2$ & $\tfrac{5}{2}$ & $1$ & $-4$ & $-4$ \\
$([X_1,X_2],[X_3,X_4])$ & $(v_1,v_4)(v_2,v_3)$ & $-\tfrac{5}{6}$ & $0$ & $\tfrac{16}{3}$ & $0$ & $\tfrac{13}{2}$ & $4$ & $-5$ & $-2$ & $8$ & $8$ \\
$([X_1,X_4],[X_2,X_3])$ & $(v_1,v_2)(v_3,v_4)$ & $\tfrac{5}{6}$ & $0$ & $-\tfrac{16}{3}$ & $0$ & $-\tfrac{13}{2}$ & $-4$ & $5$ & $2$ & $-8$ & $-8$ \\
$([X_1,X_4],[X_2,X_3])$ & $(v_1,v_3)(v_2,v_4)$ & $-\tfrac{5}{12}$ & $0$ & $\tfrac{8}{3}$ & $0$ & $\tfrac{13}{4}$ & $2$ & $-\tfrac{5}{2}$ & $-1$ & $4$ & $4$ \\
$([X_1,X_4],[X_2,X_3])$ & $(v_1,v_4)(v_2,v_3)$ & $-\tfrac{5}{12}$ & $0$ & $\tfrac{8}{3}$ & $0$ & $\tfrac{13}{4}$ & $2$ & $-\tfrac{5}{2}$ & $-1$ & $4$ & $4$ \\
averaged trace & $(v_1,v_2)(v_3,v_4)$ & $-\tfrac{37}{4}$ & $-12$ & $0$ & $\tfrac{1}{4}$ & $15$ & $24$ & $-6$ & $-12$ & $0$ & $0$ \\
averaged trace & $(v_1,v_3)(v_2,v_4)$ & $-\tfrac{37}{4}$ & $-12$ & $0$ & $\tfrac{1}{4}$ & $15$ & $24$ & $-6$ & $-12$ & $0$ & $0$ \\
averaged trace & $(v_1,v_4)(v_2,v_3)$ & $-\tfrac{37}{4}$ & $-12$ & $0$ & $\tfrac{1}{4}$ & $15$ & $24$ & $-6$ & $-12$ & $0$ & $0$ \\
\end{tabular}
\tabcaption{The $\log\varepsilon$ coefficients of the graphs with external $(A|_V)^4$ in the basis of colour and Lorentz factors, summed over each class, with the total and $8$ times the tree vertex \eqref{eq:c4-tree4}. Rows with a bracket pairing are in units of $\hbar\,\operatorname{tr}_{\mathfrak g}(\mathrm{ad}_X\mathrm{ad}_Y)/(16\pi^2(X,Y))$, rows with the averaged trace in units of $\hbar/16\pi^2$.}

\fig{diagrams/composite/quartic_AV2_AWperp2_boxes.png}{Box graphs with external $(A|_V)^2(A|_{W_{\mathbb R}^\perp})^2$; (g)--(i): gaugino loops in the three cyclic orders of the legs.}

\fig{diagrams/composite/quartic_AV2_AWperp2_triangles.png}{Triangle graphs (one quartic vertex) with external $(A|_V)^2(A|_{W_{\mathbb R}^\perp})^2$.}

\fig{diagrams/composite/quartic_AV2_AWperp2_bubbles.png}{Bubble graphs (two quartic vertices) with external $(A|_V)^2(A|_{W_{\mathbb R}^\perp})^2$.}

\begin{tabular}{l|l|l|c|c|c|c|c}
graphs & vertices & lines & number & $1/|\operatorname{Aut}\gamma|$ & $([X_1,X_2],[X_3,X_4])$ $\times\,(v_1,v_2)(v_3,v_4)$ & $([X_1,X_4],[X_2,X_3])$ $\times\,(v_1,v_2)(v_3,v_4)$ & averaged trace $\times\,(v_1,v_2)(v_3,v_4)$ \\
\hline
box (a), (b) & $(A|_V)^3$, $(A|_V)^3$, $A|_V(A|_{W_{\mathbb R}^\perp})^2$, $A|_V(A|_{W_{\mathbb R}^\perp})^2$ & $A|_V$, $A|_V$, $A|_V$, $A|_{W_{\mathbb R}^\perp}$ & $2$ & $1$ & $\tfrac{1}{8}$ & $\tfrac{1}{4}$ & $-\tfrac{3}{2}$ \\
box (c), (f) & $(A|_V)^3$, $A|_V(A|_{W_{\mathbb R}^\perp})^2$, $A|_V(A|_{W_{\mathbb R}^\perp})^2$, $A|_V(A|_{W_{\mathbb R}^\perp})^2$ & $A|_V$, $A|_V$, $A|_{W_{\mathbb R}^\perp}$, $A|_{W_{\mathbb R}^\perp}$ & $2$ & $1$ & $0$ & $0$ & $0$ \\
box (d), (e) & $A|_V(A|_{W_{\mathbb R}^\perp})^2$, $A|_V(A|_{W_{\mathbb R}^\perp})^2$, $A|_V(A|_{W_{\mathbb R}^\perp})^2$, $A|_V(A|_{W_{\mathbb R}^\perp})^2$ & $A|_V$, $A|_{W_{\mathbb R}^\perp}$, $A|_{W_{\mathbb R}^\perp}$, $A|_{W_{\mathbb R}^\perp}$ & $2$ & $1$ & $\tfrac{1}{6}$ & $\tfrac{1}{3}$ & $-2$ \\
box (g)--(i) & $\lambda\lambda A|_V$, $\lambda\lambda A|_V$, $\lambda\lambda A|_{W_{\mathbb R}^\perp}$, $\lambda\lambda A|_{W_{\mathbb R}^\perp}$ & $\lambda$, $\lambda$, $\lambda$, $\lambda$ & $3$ & $1$ & $-4$ & $-8$ & $0$ \\
triangle (a) & $(A|_V)^3$, $(A|_V)^3$, $(A|_V)^2(A|_{W_{\mathbb R}^\perp})^2$ & $A|_V$, $A|_V$, $A|_V$ & $1$ & $1$ & $-\tfrac{3}{4}$ & $-\tfrac{3}{2}$ & $9$ \\
triangle (b) & $A|_V(A|_{W_{\mathbb R}^\perp})^2$, $A|_V(A|_{W_{\mathbb R}^\perp})^2$, $(A|_{W_{\mathbb R}^\perp})^4$ & $A|_{W_{\mathbb R}^\perp}$, $A|_{W_{\mathbb R}^\perp}$, $A|_{W_{\mathbb R}^\perp}$ & $1$ & $1$ & $-\tfrac{5}{6}$ & $-\tfrac{5}{3}$ & $10$ \\
triangle (c) & $A|_V(A|_{W_{\mathbb R}^\perp})^2$, $A|_V(A|_{W_{\mathbb R}^\perp})^2$, $(A|_V)^4$ & $A|_V$, $A|_V$, $A|_{W_{\mathbb R}^\perp}$ & $1$ & $1$ & $-\tfrac{1}{8}$ & $-\tfrac{1}{4}$ & $\tfrac{3}{2}$ \\
triangle (d), (i)--(l) & $A|_V(A|_{W_{\mathbb R}^\perp})^2$, $A|_V(A|_{W_{\mathbb R}^\perp})^2$, $(A|_V)^2(A|_{W_{\mathbb R}^\perp})^2$ & $A|_V$, $A|_{W_{\mathbb R}^\perp}$, $A|_{W_{\mathbb R}^\perp}$ & $5$ & $1$ & $-\tfrac{5}{6}$ & $-\tfrac{5}{3}$ & $4$ \\
triangle (e)--(h) & $(A|_V)^3$, $A|_V(A|_{W_{\mathbb R}^\perp})^2$, $(A|_V)^2(A|_{W_{\mathbb R}^\perp})^2$ & $A|_V$, $A|_V$, $A|_{W_{\mathbb R}^\perp}$ & $4$ & $1$ & $-\tfrac{5}{4}$ & $-\tfrac{5}{2}$ & $-3$ \\
bubble (a) & $(A|_V)^4$, $(A|_V)^2(A|_{W_{\mathbb R}^\perp})^2$ & $A|_V$, $A|_V$ & $1$ & $\tfrac12$ & $\tfrac{1}{2}$ & $1$ & $-6$ \\
bubble (b) & $(A|_V)^2(A|_{W_{\mathbb R}^\perp})^2$, $(A|_{W_{\mathbb R}^\perp})^4$ & $A|_{W_{\mathbb R}^\perp}$, $A|_{W_{\mathbb R}^\perp}$ & $1$ & $\tfrac12$ & $\tfrac{5}{6}$ & $\tfrac{5}{3}$ & $-10$ \\
bubble (c), (d) & $(A|_V)^2(A|_{W_{\mathbb R}^\perp})^2$, $(A|_V)^2(A|_{W_{\mathbb R}^\perp})^2$ & $A|_V$, $A|_{W_{\mathbb R}^\perp}$ & $2$ & $1$ & $\tfrac{13}{6}$ & $\tfrac{13}{3}$ & $-2$ \\
total & & & & & $-4$ & $-8$ & $0$ \\
$8\times$tree & & & & & $-4$ & $-8$ & $0$ \\
\end{tabular}
\tabcaption{The $\log\varepsilon$ coefficients of the graphs with external $(A|_V)^2(A|_{W_{\mathbb R}^\perp})^2$ (legs $1,2$ along $V$, legs $3,4$ along $W_{\mathbb R}^\perp$), summed over each class, with the total and $8$ times the tree vertex; units as in the previous table.}

\fig{diagrams/composite/quartic_AWperp4_boxes.png}{Box graphs with external $(A|_{W_{\mathbb R}^\perp})^4$: alternating $A|_V$ and $A|_{W_{\mathbb R}^\perp}$ lines (a)--(f), gaugino loops (g)--(i).}

\fig{diagrams/composite/quartic_AWperp4_triangles.png}{Triangle graphs (one quartic vertex) with external $(A|_{W_{\mathbb R}^\perp})^4$: quartic vertex $(A|_V)^2(A|_{W_{\mathbb R}^\perp})^2$ (a)--(f) and $(A|_{W_{\mathbb R}^\perp})^4$ (g)--(l).}

\fig{diagrams/composite/quartic_AWperp4_bubbles.png}{Bubble graphs (two quartic vertices) with external $(A|_{W_{\mathbb R}^\perp})^4$: $A|_V$ lines (a)--(c), $A|_{W_{\mathbb R}^\perp}$ lines (d)--(f).}

\begin{tabular}{l|l|l|c|c}
graphs & vertices & lines & number & $1/|\operatorname{Aut}\gamma|$ \\
\hline
box (a)--(f) & $A|_V(A|_{W_{\mathbb R}^\perp})^2$, $A|_V(A|_{W_{\mathbb R}^\perp})^2$, $A|_V(A|_{W_{\mathbb R}^\perp})^2$, $A|_V(A|_{W_{\mathbb R}^\perp})^2$ & $A|_V$, $A|_V$, $A|_{W_{\mathbb R}^\perp}$, $A|_{W_{\mathbb R}^\perp}$ & $6$ & $1$ \\
box (g)--(i) & $\lambda\lambda A|_{W_{\mathbb R}^\perp}$, $\lambda\lambda A|_{W_{\mathbb R}^\perp}$, $\lambda\lambda A|_{W_{\mathbb R}^\perp}$, $\lambda\lambda A|_{W_{\mathbb R}^\perp}$ & $\lambda$, $\lambda$, $\lambda$, $\lambda$ & $3$ & $1$ \\
triangle (a)--(f) & $A|_V(A|_{W_{\mathbb R}^\perp})^2$, $A|_V(A|_{W_{\mathbb R}^\perp})^2$, $(A|_V)^2(A|_{W_{\mathbb R}^\perp})^2$ & $A|_V$, $A|_V$, $A|_{W_{\mathbb R}^\perp}$ & $6$ & $1$ \\
triangle (g)--(l) & $A|_V(A|_{W_{\mathbb R}^\perp})^2$, $A|_V(A|_{W_{\mathbb R}^\perp})^2$, $(A|_{W_{\mathbb R}^\perp})^4$ & $A|_V$, $A|_{W_{\mathbb R}^\perp}$, $A|_{W_{\mathbb R}^\perp}$ & $6$ & $1$ \\
bubble (a)--(c) & $(A|_V)^2(A|_{W_{\mathbb R}^\perp})^2$, $(A|_V)^2(A|_{W_{\mathbb R}^\perp})^2$ & $A|_V$, $A|_V$ & $3$ & $\tfrac12$ \\
bubble (d)--(f) & $(A|_{W_{\mathbb R}^\perp})^4$, $(A|_{W_{\mathbb R}^\perp})^4$ & $A|_{W_{\mathbb R}^\perp}$, $A|_{W_{\mathbb R}^\perp}$ & $3$ & $\tfrac12$ \\
\end{tabular}
\tabcaption{The classes of one-loop graphs with external $(A|_{W_{\mathbb R}^\perp})^4$.}

\begin{tabular}{l|l|c|c|c|c|c|c|c|c}
colour & Lorentz & box (a)--(f) & box (g)--(i) & triangle (a)--(f) & triangle (g)--(l) & bubble (a)--(c) & bubble (d)--(f) & total & $8\times$tree \\
\hline
$([X_1,X_2],[X_3,X_4])$ & $(v_1,v_2)(v_3,v_4)$ & $\tfrac{1}{6}$ & $-\tfrac{16}{3}$ & $-\tfrac{1}{3}$ & $-2$ & $\tfrac{2}{3}$ & $\tfrac{17}{6}$ & $-4$ & $-4$ \\
$([X_1,X_2],[X_3,X_4])$ & $(v_1,v_3)(v_2,v_4)$ & $\tfrac{1}{6}$ & $-\tfrac{16}{3}$ & $-\tfrac{1}{3}$ & $-2$ & $\tfrac{2}{3}$ & $\tfrac{17}{6}$ & $-4$ & $-4$ \\
$([X_1,X_2],[X_3,X_4])$ & $(v_1,v_4)(v_2,v_3)$ & $-\tfrac{1}{3}$ & $\tfrac{32}{3}$ & $\tfrac{2}{3}$ & $4$ & $-\tfrac{4}{3}$ & $-\tfrac{17}{3}$ & $8$ & $8$ \\
$([X_1,X_4],[X_2,X_3])$ & $(v_1,v_2)(v_3,v_4)$ & $\tfrac{1}{3}$ & $-\tfrac{32}{3}$ & $-\tfrac{2}{3}$ & $-4$ & $\tfrac{4}{3}$ & $\tfrac{17}{3}$ & $-8$ & $-8$ \\
$([X_1,X_4],[X_2,X_3])$ & $(v_1,v_3)(v_2,v_4)$ & $-\tfrac{1}{6}$ & $\tfrac{16}{3}$ & $\tfrac{1}{3}$ & $2$ & $-\tfrac{2}{3}$ & $-\tfrac{17}{6}$ & $4$ & $4$ \\
$([X_1,X_4],[X_2,X_3])$ & $(v_1,v_4)(v_2,v_3)$ & $-\tfrac{1}{6}$ & $\tfrac{16}{3}$ & $\tfrac{1}{3}$ & $2$ & $-\tfrac{2}{3}$ & $-\tfrac{17}{6}$ & $4$ & $4$ \\
averaged trace & $(v_1,v_2)(v_3,v_4)$ & $-2$ & $16$ & $4$ & $0$ & $-8$ & $-10$ & $0$ & $0$ \\
averaged trace & $(v_1,v_3)(v_2,v_4)$ & $-2$ & $16$ & $4$ & $0$ & $-8$ & $-10$ & $0$ & $0$ \\
averaged trace & $(v_1,v_4)(v_2,v_3)$ & $-2$ & $16$ & $4$ & $0$ & $-8$ & $-10$ & $0$ & $0$ \\
\end{tabular}
\tabcaption{The $\log\varepsilon$ coefficients of the graphs with external $(A|_{W_{\mathbb R}^\perp})^4$ in the basis of colour and Lorentz factors, with the total and $8$ times the tree vertex \eqref{eq:c4-tree4}; units as in the $(A|_V)^4$ table.}

In each quartic set the averaged traces cancel in the total, so that the total is a combination of the two bracket pairings, and the total equals $8$ times the tree vertex.

\subsection{The logarithmic part as a rescaling of the fields}

Write $\int(A,A^*)$, $\int(\lambda^*,\lambda)$, $\int(c^*,c)$ for the local functionals $e\mapsto\langle e_A,e_{A^*}\rangle$, $\langle e_{\lambda^*},e_\lambda\rangle$, $\langle e_{c^*},e_c\rangle$, the pairing of Definition \ref{def:c1-pairing} applied to the components of the field $e$ in the two slots of a pair, the even member first; the first uses all ten components of $A$. For a monomial in the fields, let $n_A,n_{A^*},n_\lambda,n_{\lambda^*},n_c,n_{c^*}$ denote the numbers of its factors of each kind.

\begin{lemma}[The bracket as a field count]
\label{lem:c4-count}
For every monomial $F$ in the fields,
\[
\Big\{\int(A,A^*),F\Big\}=(n_A-n_{A^*})F,\qquad \Big\{\int(\lambda^*,\lambda),F\Big\}=(n_{\lambda^*}-n_\lambda)F,\qquad \Big\{\int(c^*,c),F\Big\}=(n_{c^*}-n_c)F .
\]
Hence, on every monomial of $S$,
\[
\Big\{2\int(A,A^*)-4\int(\lambda^*,\lambda)-\int(c^*,c),\,F\Big\}=\big(2(n_A-n_{A^*})+4(n_\lambda-n_{\lambda^*})+(n_c-n_{c^*})\big)F .
\]
\end{lemma}

\begin{proof}
By Definition \ref{def:c1-bracket}, the kernel $K_0$ pairs each even member $A,\lambda^*,c^*$ with its odd partner $A^*,\lambda,c$ with the coefficient $-\delta(x-y)$ times the inverse fibre pairing, which is the structure $K_0=-\sum(x_i\otimes\xi_i+\xi_i\otimes x_i)$, $\Delta_0=\sum\partial_{x_i}\partial_{\xi_i}$ of [C07, §10.1] for the manifold a point, in the notation of loc. cit.: $x_i$ an even coordinate and $\xi_i$ its odd partner, applied pointwise at each position. For one pair, $x_i\xi_i$ is odd, and
\[
\begin{aligned}
\{x_i\xi_i,x_i\}&=\Delta_0(x_i^2\xi_i)-\Delta_0(x_i\xi_i)\,x_i+x_i\xi_i\,\Delta_0(x_i)=2x_i-x_i=x_i,\\
\{x_i\xi_i,\xi_i\}&=\Delta_0(x_i\xi_i\xi_i)-\Delta_0(x_i\xi_i)\,\xi_i+x_i\xi_i\,\Delta_0(\xi_i)=-\xi_i .
\end{aligned}
\]
The bracket with a fixed functional is a derivation, because $\Delta_0$ is a second-order operator, so $\{x_i\xi_i,-\}=x_i\partial_{x_i}-\xi_i\partial_{\xi_i}$ counts $x_i$ minus $\xi_i$. The three functionals of the lemma are the sums of $x_i\xi_i$ over the components, with $x_i$ the even member, which gives the three identities; the combination follows by linearity.
\end{proof}

\begin{theorem}[The logarithmic part of the one-loop counterterm]
\label{thm:c4-log}
The $\log\varepsilon$ part of $I^{CT}[\varepsilon]$ is
\begin{equation}
\frac{\hbar\,\operatorname{tr}_{\mathfrak g}(\mathrm{ad}_X\mathrm{ad}_Y)}{16\pi^2(X,Y)}\,\log\varepsilon\,\big(Q+\{I,-\}\big)\Big(2\int(A,A^*)-4\int(\lambda^*,\lambda)-\int(c^*,c)\Big).
\label{eq:c4-log}
\end{equation}
Equivalently, it multiplies each monomial of $S$ by $2(n_A-n_{A^*})+4(n_\lambda-n_{\lambda^*})+(n_c-n_{c^*})$: the bracket acts as a weighted field count with weights $2$ for $A$, $4$ for $\lambda$, $1$ for $c$, and the opposite weights for the antifields.
\end{theorem}

\begin{proof}
The functional $2\int(A,A^*)-4\int(\lambda^*,\lambda)-\int(c^*,c)$ is odd. By Proposition \ref{prop:c1-Q}, $Q$ applied to it equals its bracket with $\tfrac12\langle e,Qe\rangle$, so
\[
\begin{aligned}
\big(Q+\{I,-\}\big)\Big(2\int(A,A^*)-4\int(\lambda^*,\lambda)-\int(c^*,c)\Big)&=\Big\{S,\,2\int(A,A^*)-4\int(\lambda^*,\lambda)-\int(c^*,c)\Big\}\\
&=\Big\{2\int(A,A^*)-4\int(\lambda^*,\lambda)-\int(c^*,c),\,S\Big\}
\end{aligned}
\]
by the graded symmetry of Definition \ref{def:c1-bracket}, $S$ being even. By Lemma \ref{lem:c4-count}, this bracket multiplies the terms of $S$ by
\[
\begin{aligned}
&-\tfrac14\int\mathrm{dvol}\,(dA,dA):\ 4,\qquad \tfrac12\int\mathrm{dvol}\,(\lambda,{\not d}\lambda):\ 8,\qquad \int\mathrm{dvol}\,(dc,A^*):\ 1-2=-1,\\
&\int\mathrm{dvol}\,\big(-\tfrac12(dA,\tfrac12[A\wedge A])\big):\ 6,\qquad \int\mathrm{dvol}\,\tfrac12(\lambda,({\not d}_A-{\not d})\lambda):\ 8+2=10,\\
&\int\mathrm{dvol}\,([A,c],A^*):\ 2+1-2=1,\qquad \int\mathrm{dvol}\,([c,\lambda],\lambda^*):\ 1+4-4=1,\\
&\int\mathrm{dvol}\,\tfrac12([c,c],c^*):\ 2-1=1,\qquad \int\mathrm{dvol}\,\big(-\tfrac14(\tfrac12[A\wedge A],\tfrac12[A\wedge A])\big):\ 8 .
\end{aligned}
\]
These are the computed multiples: the quadratic ones by Theorem \ref{thm:c4-quadratic} and Remark \ref{rem:c4-quadratic} (with the coefficient $0$ of $\int\mathrm{dvol}\,(d^*A|_V,d^*A|_V)$, absent from $S$), the cubic ones $6,6,10,10,1,1,1,1$ by the totals of the eight cubic tables, and the quartic ones $8,8,8$ by the totals of the three quartic tables. In the cubic and quartic sets the computed coefficient equals the multiple of the tree vertex as a multilinear form, including the two chirality structures of $(\lambda,\lambda,A|_{W_{\mathbb R}^\perp})$ and of $(c,\lambda,\lambda^*)$. By Proposition \ref{prop:c4-structure} these leg sets carry the whole $\log\varepsilon$ part of $I^{CT}[\varepsilon]$, so the two local functionals coincide.
\end{proof}

\begin{corollary}[The coupling counterterm vanishes at one loop]
\label{cor:c4-coupling}
Suppose the $\log\varepsilon$ part of $I^{CT}[\varepsilon]$ is written as $\frac{\hbar\,\operatorname{tr}_{\mathfrak g}(\mathrm{ad}_X\mathrm{ad}_Y)}{16\pi^2(X,Y)}\log\varepsilon$ times the sum of $\{S,-\}$ of an element of the span of $\int(A|_V,A^*|_V)$, $\int(A|_{W_{\mathbb R}^\perp},A^*|_{W_{\mathbb R}^\perp})$, $\int(c^*,c)$ and the two chiral parts of $\int(\lambda^*,\lambda)$, and a multiple of $S$. Then the multiple of $S$ is $0$ and the element is $2\int(A,A^*)-4\int(\lambda^*,\lambda)-\int(c^*,c)$. The same holds with $S$ replaced by its classical part, the terms of $S$ without antifields. Since the coupling constant sits in $(-,-)$ and enters $S$ as an overall factor, the one-loop $\log\varepsilon$ counterterm is a rescaling of the fields with vanishing component along the coupling.
\end{corollary}

\begin{proof}
By Lemma \ref{lem:c4-count}, $\{S,-\}$ of an element of the span multiplies each monomial of $S$ by a sum of field counts weighted by the coefficients of the element, additive over the factors of the monomial, and a multiple of $S$ adds the same number to every monomial. The terms of $S$ built from $A|_V$ alone contain $2$, $3$, $4$ factors $A|_V$, with computed multiples $4$, $6$, $8$: an additive weight times $2,3,4$ plus a common constant fits $4,6,8$ with weight $2$ and constant $3\cdot4-2\cdot6=0$. With the constant $0$: the kinetic terms of $A|_{W_{\mathbb R}^\perp}$ and $\lambda$ give the weight $2$ of $A|_{W_{\mathbb R}^\perp}$ and the sum $8$ of the two chiral weights of $\lambda$; the vertex $\lambda\lambda A|_{W_{\mathbb R}^\perp}$, whose two structures have both $\lambda$ of the same chirality, gives each chiral weight $(10-2)/2=4$; the ghost kinetic term $\int\mathrm{dvol}\,(dc,A^*)$ gives the weight $-1+2=1$ of $c$; the remaining eight multiples are then reproduced. For the classical part: by Lemma \ref{lem:c4-count} every term of $S$ containing an antifield has $n_c-n_{c^*}=1$ and every classical term has $n_c=n_{c^*}=0$, so $S$ minus its classical part equals $\{S,-\int(c^*,c)\}$, an element of the span; replacing $S$ by its classical part shifts the element by a multiple of $\int(c^*,c)$ and leaves the multiple unchanged.
\end{proof}

\begin{remark}[The mass term has nonzero BV variation]
\label{rem:c4-epsinv}
The $\varepsilon^{-1}$ part of $I^{CT}[\varepsilon]$, the mass term $\frac{\hbar\,\operatorname{tr}_{\mathfrak g}(\mathrm{ad}_X\mathrm{ad}_Y)}{16\pi^2(X,Y)}\tfrac94\varepsilon^{-1}\int\mathrm{dvol}\,(A,A)$, has nonzero variation under $Q+\{I,-\}$. Indeed $\{I,\int\mathrm{dvol}\,(A,A)\}$ pairs the $A^*$ of $\int\mathrm{dvol}\,([A,c],A^*)$ with $A$ and is a multiple of $\int\mathrm{dvol}\,([A,c],A)=-\int\mathrm{dvol}\,(c,[A,A])=0$, the bracket $[A,A]$ contracting the symmetric metric with the antisymmetric bracket of $\mathfrak g$. By Proposition \ref{prop:c1-Q}, $Q$ inserts $Qc=-dc$ into the slot of $A|_V$, so
\[
\big(Q+\{I,-\}\big)\Big(\frac{\hbar\,\operatorname{tr}_{\mathfrak g}(\mathrm{ad}_X\mathrm{ad}_Y)}{16\pi^2(X,Y)}\tfrac94\varepsilon^{-1}\int\mathrm{dvol}\,(A,A)\Big)=-\frac{\hbar\,\operatorname{tr}_{\mathfrak g}(\mathrm{ad}_X\mathrm{ad}_Y)}{16\pi^2(X,Y)}\tfrac92\varepsilon^{-1}\int\mathrm{dvol}\,(dc,A|_V).
\]
This is the breaking of the gauge symmetry by the regulator. The next section shows how the master equation at the UV boundary absorbs it and what it leaves behind.
\end{remark}

\section{The quantum master equation at the UV boundary}

\begin{definition}[Defect vertex]
\label{def:c4-defect}
For an even functional $J$ with values in $\mathbb C[[\hbar]]$, at least cubic modulo $\hbar$, the defect vertex at scale $\varepsilon$ is the odd functional
\[
QJ+\tfrac12\{J,J\}_\varepsilon+\hbar\Delta_\varepsilon J,
\]
with $\Delta_\varepsilon$ and $\{-,-\}_\varepsilon$ of Definition \ref{def:c3-delta}. By Corollary \ref{cor:c3-qme}, its vanishing is the quantum master equation at scale $\varepsilon$.
\end{definition}

\begin{proposition}[The $Q$-variation at scale $\infty$ is the flow of the defect]
\label{prop:c4-defectflow}
On functionals without superconformal ghosts,
\begin{equation}
Q\,W(P(\varepsilon,\infty),J)=e^{-W(P(\varepsilon,\infty),J)/\hbar}\,e^{\hbar\partial_{P(\varepsilon,\infty)}}\Big(\big(QJ+\tfrac12\{J,J\}_\varepsilon+\hbar\Delta_\varepsilon J\big)\,e^{J/\hbar}\Big).
\label{eq:c4-defectflow}
\end{equation}
Thus $Q\,W(P(\varepsilon,\infty),J)$ is the sum of the connected graphs with exactly one vertex labelled by the defect vertex, all other vertices by $J$, and all edges by $P(\varepsilon,\infty)$, with the weights of Proposition \ref{prop:c3-graphs} [C07, §15.4].
\end{proposition}

\begin{proof}
For an even $J$, the product rule of Definition \ref{def:c3-delta} gives $(Q+\hbar\Delta_\varepsilon)e^{J/\hbar}=\hbar^{-1}\big(QJ+\tfrac12\{J,J\}_\varepsilon+\hbar\Delta_\varepsilon J\big)e^{J/\hbar}$, as in the proof of Corollary \ref{cor:c3-qme}. Theorem \ref{thm:c3-uv} without superconformal ghosts gives
\[
Q\,e^{\hbar\partial_{P(\varepsilon,\infty)}}e^{J/\hbar}=e^{\hbar\partial_{P(\varepsilon,\infty)}}\big(Q+\hbar\Delta_\varepsilon\big)e^{J/\hbar}=\hbar^{-1}e^{\hbar\partial_{P(\varepsilon,\infty)}}\Big(\big(QJ+\tfrac12\{J,J\}_\varepsilon+\hbar\Delta_\varepsilon J\big)\,e^{J/\hbar}\Big).
\]
The left side is $Q\,e^{W(P(\varepsilon,\infty),J)/\hbar}=\hbar^{-1}\big(Q\,W(P(\varepsilon,\infty),J)\big)e^{W(P(\varepsilon,\infty),J)/\hbar}$, $Q$ being a derivation; multiplying by $\hbar\,e^{-W(P(\varepsilon,\infty),J)/\hbar}$ gives \eqref{eq:c4-defectflow}. Expanding $e^{\hbar\partial_{P(\varepsilon,\infty)}}$ into graphs as in Proposition \ref{prop:c3-graphs}, the right side of \eqref{eq:c4-defectflow} is the sum of the graphs with exactly one vertex labelled by the defect vertex; the factor $e^{-W(P(\varepsilon,\infty),J)/\hbar}$, the exponential of the connected graphs of $J$, removes the components without that vertex, which leaves the connected graphs.
\end{proof}

With $\Delta_\infty=0$ and $\{-,-\}_\infty=0$, the quantum master equation at scale $\infty$ for the renormalised flow is $Q\lim_{\varepsilon\to0}W(P(\varepsilon,\infty),J)=0$ (Corollary \ref{cor:c3-qme}). Its failure at order $\hbar$ is the order-$\hbar$ part of $\lim_{\varepsilon\to0}Q\,W(P(\varepsilon,\infty),I-I^{CT}[\varepsilon])$. By \eqref{eq:c4-defectflow} it is a sum of graphs with one defect vertex, among them trees of $I$ attached to that vertex; it is the flow to scale $\infty$ of an effective defect vertex, the one-loop obstruction. On a compact manifold with a finite upper scale the obstruction is the local functional of [C07, Prop 12.0.3], closed and independent of $\varepsilon$ by [C07, Lemma 12.0.4]; on $\mathbb R^4$ with upper scale $\infty$ these properties are part of Step 2. In the sector with legs $(c,A|_V)$ no tree of $I$ contributes at order $\hbar$, and the limit computed below is itself local.

\begin{proposition}[The defect vertex at one loop]
\label{prop:c4-defect1}
For $J=I-I^{CT}[\varepsilon]$, with $I$ the classical interaction \eqref{eq:c1-I},
\begin{equation}
QJ+\tfrac12\{J,J\}_\varepsilon+\hbar\Delta_\varepsilon J=\tfrac12\big(\{I,I\}_\varepsilon-\{I,I\}\big)+\hbar\Delta_\varepsilon I-Q\,I^{CT}[\varepsilon]-\{I,I^{CT}[\varepsilon]\}_\varepsilon+(\text{terms of order }\hbar^2).
\label{eq:c4-defect1}
\end{equation}
The first term is the bracket built from $\Delta_\varepsilon-\Delta_0=-\partial_{K_\varepsilon-K_0}$, i.e. two vertices of $I$ joined by one special edge carrying
\[
K_\varepsilon-K_0=-\int_0^\varepsilon(H\otimes1)K_t\,dt,
\]
which in momentum space is $e^{-\varepsilon|p|^2}-1$ times the components of $K_0$; its Schwinger time lies in $[0,\varepsilon]$.
\end{proposition}

\begin{proof}
Expanding $QJ+\tfrac12\{J,J\}_\varepsilon+\hbar\Delta_\varepsilon J$ with $J=I-I^{CT}[\varepsilon]$ and $I^{CT}[\varepsilon]$ of order $\hbar$: at order $\hbar^0$, $QI+\tfrac12\{I,I\}_\varepsilon$; at order $\hbar$, $-QI^{CT}[\varepsilon]-\tfrac12\big(\{I,I^{CT}[\varepsilon]\}_\varepsilon+\{I^{CT}[\varepsilon],I\}_\varepsilon\big)+\hbar\Delta_\varepsilon I$, where the two brackets of even functionals coincide; the rest is $\tfrac12\{I^{CT}[\varepsilon],I^{CT}[\varepsilon]\}_\varepsilon-\hbar\Delta_\varepsilon I^{CT}[\varepsilon]$, of order $\hbar^2$. The classical master equation (Proposition \ref{prop:c1-cme}) gives $QI=-\tfrac12\{I,I\}$, with $\{-,-\}$ built from $\Delta_0$ (Definition \ref{def:c1-bracket}). From $K_t=(e^{-tH}\otimes1)K_0$, $\frac{d}{dt}K_t=-(H\otimes1)K_t$, and integration from $0$ to $\varepsilon$ gives the special edge; in momentum space $e^{-tH}$ is $e^{-t|p|^2}$.
\end{proof}

At order $\hbar$, the graphs of \eqref{eq:c4-defectflow} with two legs are: the term $-QI^{CT}[\varepsilon]$ as a vertex; the bubbles formed by the two vertices of $\tfrac12(\{I,I\}_\varepsilon-\{I,I\})$, joined by the special edge and one edge $P(\varepsilon,\infty)$; and one-particle-reducible graphs in which a loop closes at a single cubic vertex. The term $\{I,I^{CT}[\varepsilon]\}_\varepsilon$ has at least three legs.

\begin{lemma}[Two-leg sectors of degree one]
\label{lem:c4-sectors}
An invariant local functional of degree $1$ with two legs and weight at most $4$ is a combination of $\int\mathrm{dvol}\,(dc,A|_V)$ (weight $2$) and $\int\mathrm{dvol}\,(d^*dc,d^*A|_V)$ (weight $4$). In particular the sectors with legs $(c,A|_{W_{\mathbb R}^\perp})$ and $(c,\lambda)$ vanish.
\end{lemma}

\begin{proof}
Degree $1$ with two legs requires one $c$ and one leg of degree $0$, $A$ or $\lambda$. The invariant pairings of a scalar with a spinor and with a vector of $W_{\mathbb R}^\perp$ vanish. With $A|_V$: one derivative in weight $2$, $\int\mathrm{dvol}\,(dc,A|_V)$; three derivatives in weight $4$, where $\int\mathrm{dvol}\,(dc,d^*dA|_V)=\int\mathrm{dvol}\,(d\,dc,dA|_V)=0$ and $\int\mathrm{dvol}\,(dc,dd^*A|_V)=\int\mathrm{dvol}\,(d^*dc,d^*A|_V)$.
\end{proof}

\fig{diagrams/composite/qme_defect_c_AV.png}{The one-loop graphs of the defect vertex with external $c,A|_V$: the shaded edge is the special edge $K_\varepsilon-K_0$ of $\tfrac12(\{I,I\}_\varepsilon-\{I,I\})$, the other edge carries $P(\varepsilon,\infty)$. Special edge between the slots (a) $A|_V$, $A^*|_V$ with an $A|_V$ line; (b) $A|_V$, $A^*|_V$ with a $c$--$A^*$ line; (c) $A|_{W_{\mathbb R}^\perp}$, $A^*|_{W_{\mathbb R}^\perp}$ with an $A|_{W_{\mathbb R}^\perp}$ line; (d) $\lambda^*$, $\lambda$ with a $\lambda$ line; (e) $c^*$, $c$ with a $c$--$A^*$ line.}

\fig{diagrams/composite/qme_defect_c_AV_reducible.png}{The one-particle-reducible graphs of the defect vertex with external $c,A|_V$. Row 1: from $\tfrac12(\{I,I\}_\varepsilon-\{I,I\})$, with an edge $P(\varepsilon,\infty)$ closing at one cubic vertex. Row 2: from $\hbar\Delta_\varepsilon I$, a special edge $K_\varepsilon$ closing at one cubic vertex. Each closed loop at a cubic vertex carries the colour factor $\operatorname{tr}_{\mathfrak g}(\mathrm{ad}_X)=0$.}

\begin{lemma}[Graphs with one special edge]
\label{lem:c4-special}
The plane-wave value of each bubble of the first figure is $\varepsilon^{-1}$ times a multiple of $\int\mathrm{dvol}\,(dc,A|_V)$ plus a multiple of $\int\mathrm{dvol}\,(d^*dc,d^*A|_V)$ plus a term bounded by a multiple of $\varepsilon$, with vanishing coefficient of $\log\varepsilon$.
\end{lemma}

\begin{proof}
The special edge carries $-\int_0^\varepsilon(\cdots)e^{-t_1|p_1|^2}dt_1$ and the other edge $\int_\varepsilon^\infty(\cdots)e^{-t_2|p_1-p|^2}dt_2$. After the Gaussian integral over $p_1$ (Lemma \ref{lem:c4-gauss}, extended to the polynomial numerators of degree three of these graphs by the same moments), the integrand is $(t_1+t_2)^{-2}e^{-t_1t_2|p|^2/(t_1+t_2)}$ times a sum of terms $(t_1+t_2)^{-m}$, $m\ge0$, with polynomial coefficients in $\frac{t_1}{t_1+t_2}$ and $p$. On $[0,\varepsilon]\times[\varepsilon,\infty)$ the exponent satisfies $0\le\frac{t_1t_2|p|^2}{t_1+t_2}\le t_1|p|^2\le\varepsilon|p|^2$, so its exponential is a power series in $\varepsilon|p|^2$ with uniform remainders. Substituting $\varepsilon t_1$, $\varepsilon t_2$ for $t_1$, $t_2$ maps the region to $[0,1]\times[1,\infty)$, on which $(t_1+t_2)^{-2-m}$ is integrable, and each term becomes an integer power of $\varepsilon$ times a convergent integral. Hence the value is a series in integer powers of $\varepsilon$; by Lemma \ref{lem:c4-sectors} the coefficients of $\varepsilon^{-1}$ and $\varepsilon^0$ are the two listed functionals.
\end{proof}

\begin{tabular}{l|l|l|c|c}
graph & special edge $K_\varepsilon-K_0$ & edge $P(\varepsilon,\infty)$ & coefficient of $\varepsilon^{-1}\int\mathrm{dvol}\,(dc,A|_V)$ & coefficient of $\int\mathrm{dvol}\,(d^*dc,d^*A|_V)$ \\
\hline
(a) & $A|_V$--$A^*|_V$ & $A|_V$ & $-\tfrac94$ & $-\tfrac58$ \\
(b) & $A|_V$--$A^*|_V$ & $c$--$A^*$ & $\tfrac18$ & $-\tfrac1{12}$ \\
(c) & $A|_{W_{\mathbb R}^\perp}$--$A^*|_{W_{\mathbb R}^\perp}$ & $A|_{W_{\mathbb R}^\perp}$ & $-\tfrac92$ & $-\tfrac54$ \\
(d) & $\lambda^*$--$\lambda$ & $\lambda$ & $2$ & $-\tfrac43$ \\
(e) & $c^*$--$c$ & $c$--$A^*$ & $\tfrac18$ & $-\tfrac1{12}$ \\
$-QI^{CT}[\varepsilon]$ & & & $\tfrac92$ & $0$ \\
total & & & $0$ & $-\tfrac{27}{8}$ \\
\end{tabular}
\tabcaption{The one-loop part of $Q\,W(P(\varepsilon,\infty),I-I^{CT}[\varepsilon])$ with legs $(c,A|_V)$, graph by graph, in units of $\hbar\,\operatorname{tr}_{\mathfrak g}(\mathrm{ad}_X\mathrm{ad}_Y)/(16\pi^2(X,Y))$. Each bubble has weight $1$: the $\tfrac12$ of $\tfrac12\{I,I\}$ cancels against the two assignments of the two vertices. Every entry has vanishing $\log\varepsilon$ coefficient.}

\begin{theorem}[The one-loop obstruction in the sector $(c,A|_V)$]
\label{thm:c4-obstruction}
The part of $Q\,W(P(\varepsilon,\infty),I-I^{CT}[\varepsilon])$ of order $\hbar$ with legs $(c,A|_V)$ has the following $\varepsilon$-expansion: the special-edge bubbles contribute $-\tfrac92\,\varepsilon^{-1}$ times $\frac{\hbar\,\operatorname{tr}_{\mathfrak g}(\mathrm{ad}_X\mathrm{ad}_Y)}{16\pi^2(X,Y)}\int\mathrm{dvol}\,(dc,A|_V)$, the vertex $-QI^{CT}[\varepsilon]$ contributes $+\tfrac92\,\varepsilon^{-1}$ times the same functional, the coefficient of $\log\varepsilon$ is $0$, and the limit is
\begin{equation}
-\frac{27}{8}\,\frac{\hbar\,\operatorname{tr}_{\mathfrak g}(\mathrm{ad}_X\mathrm{ad}_Y)}{16\pi^2(X,Y)}\int\mathrm{dvol}\,(d^*dc,d^*A|_V)=\frac{27}{16}\,\frac{\hbar\,\operatorname{tr}_{\mathfrak g}(\mathrm{ad}_X\mathrm{ad}_Y)}{16\pi^2(X,Y)}\;Q\int\mathrm{dvol}\,(d^*A|_V,d^*A|_V).
\label{eq:c4-obstruction}
\end{equation}
The sectors $(c,A|_{W_{\mathbb R}^\perp})$ and $(c,\lambda)$ of the one-loop obstruction vanish.
\end{theorem}

\begin{proof}
By Proposition \ref{prop:c4-defectflow} and \eqref{eq:c4-defect1}, the order-$\hbar$ part with two legs is the sum of the graphs listed before Lemma \ref{lem:c4-sectors}. The one-particle-reducible graphs vanish: a loop closing at one cubic vertex contracts two legs of a vertex of the form $([X,Y],Z)$ with $(-,-)^{-1}$ and gives $\operatorname{tr}_{\mathfrak g}(\mathrm{ad}_X)=0$ or the contraction of the antisymmetric bracket with the symmetric $(-,-)^{-1}$, both zero. The same argument gives $\Delta_\varepsilon I=0$: $K_\varepsilon$ pairs a field with its antifield, the vertices of $I$ containing an antifield are $([A,c],A^*)$, $([c,\lambda],\lambda^*)$, $\tfrac12([c,c],c^*)$, and each self-contraction is a trace of $\mathrm{ad}$. The five bubbles are given in the table; their $\varepsilon^{-1}$ coefficients sum to $-\tfrac94+\tfrac18-\tfrac92+2+\tfrac18=-\tfrac92$, and by Lemma \ref{lem:c4-special} they have vanishing $\log\varepsilon$ coefficient. By Theorem \ref{thm:c4-quadratic} and Remark \ref{rem:c4-epsinv}, $-QI^{CT}[\varepsilon]$ restricted to $(c,A|_V)$ is $-\tfrac94\cdot(-2)=\tfrac92$ times $\varepsilon^{-1}\int\mathrm{dvol}\,(dc,A|_V)$ in the unit; its $\log\varepsilon$ part restricted to $(c,A|_V)$ vanishes, because $Q\int\mathrm{dvol}\,(dA,dA)$ contains $d\,dc=0$, $Q\int\mathrm{dvol}\,(dc,A^*)$ restricted to $(c,A|_V)$ is $-\tfrac12\int\mathrm{dvol}\,(dc,d^*dA|_V)=-\tfrac12\int\mathrm{dvol}\,(d\,dc,dA|_V)=0$, and the coefficient of $\int\mathrm{dvol}\,(d^*A|_V,d^*A|_V)$ is $0$. The finite parts sum to $-\tfrac58-\tfrac1{12}-\tfrac54-\tfrac43-\tfrac1{12}=-\tfrac{81}{24}=-\tfrac{27}8$. Finally $Q\int\mathrm{dvol}\,(d^*A|_V,d^*A|_V)=2\int\mathrm{dvol}\,(d^*A|_V,d^*(-dc))=-2\int\mathrm{dvol}\,(d^*dc,d^*A|_V)$, so $-\tfrac{27}8=\tfrac{27}{16}\cdot(-2)$. The last statement is Lemma \ref{lem:c4-sectors}.
\end{proof}

\begin{remark}[The obstruction as the longitudinal part of the two-point function]
\label{rem:c4-selfenergy}
$Q$ preserves the polynomial degree of functionals, so the part of $Q\,W(P(\varepsilon,\infty),I-I^{CT}[\varepsilon])$ with legs $(c,A|_V)$ at order $\hbar$ is $Q$ applied to the order-$\hbar$ quadratic part of $W(P(\varepsilon,\infty),I-I^{CT}[\varepsilon])$: the renormalised two-point function, whose limit as $\varepsilon\to0$ at nonzero momentum exists by Theorem \ref{thm:c4-quadratic} and Proposition \ref{prop:c4-limit}. $Q$ annihilates its $(c,A^*)$ part on $(c,A|_V)$, since $\iota_p$ of $|p|^2A|_V-p\,\iota_pA|_V$ is $0$, and its transverse $A|_V$ part, which depends on $A|_V$ through $dA|_V$. Hence the cancellation of $\varepsilon^{-1}$ and the vanishing of $\log\varepsilon$ in Theorem \ref{thm:c4-obstruction} follow from Theorem \ref{thm:c4-quadratic}, and the obstruction is $Q$ of the longitudinal part of the renormalised $A|_V$ two-point function at scale $\infty$. A longitudinal quadratic functional of $A|_V$ is $\int\mathrm{dvol}\,(d^*A|_V,G\,d^*A|_V)$ for a translation-invariant operator $G$ on functions, and $Q$ sends it to $-2\int\mathrm{dvol}\,(d^*dc,G\,d^*A|_V)$; comparing with \eqref{eq:c4-obstruction} shows that $G$ is the constant $\tfrac{27}{16}$ in the unit, so the longitudinal part is the local functional
\[
\frac{27}{16}\,\frac{\hbar\,\operatorname{tr}_{\mathfrak g}(\mathrm{ad}_X\mathrm{ad}_Y)}{16\pi^2(X,Y)}\int\mathrm{dvol}\,(d^*A|_V,d^*A|_V).
\]
Equivalently, $Q$ applied to the longitudinal parts of the graphs of the $A|_V$ two-point function, with every edge cut off at $\varepsilon$, sums to $-\tfrac92\varepsilon^{-1}$ and $-\tfrac{27}8$ times the functionals of \eqref{eq:c4-obstruction}; there the gauge bubble and the ghost loop carry separate $\log(\varepsilon|p|^2)$ and $\log2$ terms in their longitudinal parts, which cancel in their sum. Theorem \ref{thm:c4-obstruction} is the evaluation of the same quantity through the UV-boundary identity, on graphs with one special edge confined to $[0,\varepsilon]$.
\end{remark}

\begin{theorem}[The finite counterterm in the sector $(c,A|_V)$]
\label{thm:c4-finite}
Replace the renormalised interaction $I-I^{CT}[\varepsilon]$ by
\begin{equation}
I-\frac{27}{16}\,\frac{\hbar\,\operatorname{tr}_{\mathfrak g}(\mathrm{ad}_X\mathrm{ad}_Y)}{16\pi^2(X,Y)}\int\mathrm{dvol}\,(d^*A|_V,d^*A|_V)-I^{CT}[\varepsilon].
\label{eq:c4-finite}
\end{equation}
Then $I^{CT}[\varepsilon]$ at one loop is unchanged, the order-$\hbar$ part of the flow has a limit at generic leg momenta, and the part with legs $(c,A|_V)$ of the one-loop obstruction vanishes. The added term is local, even, of degree $0$ and weight $4$, independent of $\varepsilon$, and invariant under translations, $\mathrm{Spin}(4,\mathbb C)\times\mathrm{Spin}(6,\mathbb C)$, the adjoint action of $\mathfrak g$ and the dilatation. Among the quadratic functionals of $A|_V$ of weight $4$ and degree $0$, it is unique up to adding multiples of $\int\mathrm{dvol}\,(dA|_V,dA|_V)$, which is $Q$-closed.
\end{theorem}

\begin{proof}
The added term is of order $\hbar$, so at order $\hbar$ it enters the flow as a tree vertex, with a limit as $\varepsilon\to0$, and leaves the one-loop singular parts unchanged. In \eqref{eq:c4-defect1} it adds $-Q$ of itself at order $\hbar$, and $-\{I,\cdot\}_\varepsilon$ of itself, which has at least three legs. By Theorem \ref{thm:c4-obstruction} and $Q\int\mathrm{dvol}\,(d^*A|_V,d^*A|_V)=-2\int\mathrm{dvol}\,(d^*dc,d^*A|_V)$, the $(c,A|_V)$ part becomes $-\tfrac{27}{8}+\tfrac{27}{16}\cdot2=0$ in the unit. Invariance under the dilatation holds because the weight is $4$. For uniqueness: by Lemma \ref{lem:c4-enumeration}, the quadratic functionals of $A|_V$ of weight $4$ are spanned by $\int\mathrm{dvol}\,(dA|_V,dA|_V)$ and $\int\mathrm{dvol}\,(d^*A|_V,d^*A|_V)$; the first is $Q$-closed, since $d\,dc=0$, and the second has $Q$-variation $-2\int\mathrm{dvol}\,(d^*dc,d^*A|_V)\neq0$, so the coefficient $\tfrac{27}{16}$ of the second is forced.
\end{proof}

\begin{remark}[Renormalisation conditions belong to the definition of the theory]
\label{rem:c4-condition}
$I^{CT}[\varepsilon]$ is purely singular, and unique among the translation-invariant and $\mathfrak g$-invariant local functionals of Proposition \ref{prop:c4-unique} (the one-loop analogue of [C07, Thm B], in the scope of Remark \ref{rem:c4-thmB}); the finite term of \eqref{eq:c4-finite} is a finite change of the interaction, $I\mapsto I-\tfrac{27}{16}\frac{\hbar\,\operatorname{tr}_{\mathfrak g}(\mathrm{ad}_X\mathrm{ad}_Y)}{16\pi^2(X,Y)}\int\mathrm{dvol}\,(d^*A|_V,d^*A|_V)$, of the kind that lifts a theory across a $Q$-exact obstruction in [C07, Lemma 12.0.4]. The order-$\hbar$ part of the regulated flow with $I-I^{CT}[\varepsilon]$ has a finite limit at generic leg momenta (Proposition \ref{prop:c4-limit}) and fails the master equation at scale $\infty$ by a finite local amount; the master equation forces the coefficient $\tfrac{27}{16}$ and leaves the $Q$-closed finite counterterms, such as multiples of $\int\mathrm{dvol}\,(dA,dA)$, free. This choice is a renormalisation condition, and it is part of the definition of the quantum theory: every observable $O[\infty]$ is built from the flow of the interaction, and changes with it. In this program it is the first instance of that statement; the same mechanism recurs for the equivariant counterterms of the superconformal ghosts.
\end{remark}

\begin{remark}[The cubic sector of the obstruction]
\label{rem:c4-cubic}
The one-loop obstruction also has components with three legs: $(c,A|_V,A|_V)$, $(c,A|_{W_{\mathbb R}^\perp},A|_{W_{\mathbb R}^\perp})$, $(c,\lambda,\lambda)$ and $(c,c,A^*|_V)$. They are sums of triangles and quartic-vertex bubbles with one special edge, of the vertex $-\{I,I^{CT}[\varepsilon]\}_\varepsilon$, and of one-particle-reducible graphs in which $-QI^{CT}[\varepsilon]$ is joined to $I$ by an edge $P(\varepsilon,\infty)$. The finite term of \eqref{eq:c4-finite} changes them by $-\{I,\cdot\}$ of itself, a nonzero functional built from $\int\mathrm{dvol}\,([A,c],A^*)$. The full one-loop master equation, i.e. the vanishing of the obstruction in all sectors after a finite counterterm, is the subject of Step 2.
\end{remark}

\begin{proposition}[The finite counterterm under special conformal transformations]
\label{prop:c4-K}
Let $v$ be the special conformal field $x\mapsto(x,x)v_1-2(v_1,x)x$, $v_1\in V$ (Definition \ref{def:c2-even}), acting on $A|_V$ by the Lie derivative $L_v$ (Definition \ref{def:c2-S1}). The derivative of $\int\mathrm{dvol}\,(d^*A|_V,d^*A|_V)$ along $A|_V\mapsto L_vA|_V$ is
\[
2\int\mathrm{dvol}\,(d^*A|_V,d^*L_vA|_V)=-8\int\mathrm{dvol}\,(d^*A|_V,\iota_{v_1}A|_V)\neq0,
\]
while for translations, rotations and the dilatation $2E$ the derivative vanishes.
\end{proposition}

\begin{proof}
For a conformal Killing field $v$ of $V$, the symmetric part of the derivative of $(v,-)$ is $\tfrac14\operatorname{div}v$ times the metric; taking one more divergence, $d^*d$ applied to the components of $v$ is $\tfrac12\operatorname{grad}(\operatorname{div}v)$. The Lie derivative of the one-form $A|_V$ is $v(A|_V)+A|_V\circ(\text{derivative of }v)$, and the divergence of this expression, with the two identities above, is
\[
d^*(L_vA|_V)=v(d^*A|_V)+\tfrac12(\operatorname{div}v)\,d^*A|_V+\tfrac12\,\iota_{\operatorname{grad}(\operatorname{div}v)}A|_V .
\]
Hence
\[
\begin{aligned}
2\int\mathrm{dvol}\,(d^*A|_V,d^*L_vA|_V)={}&\int\mathrm{dvol}\,v\big((d^*A|_V,d^*A|_V)\big)+\int\mathrm{dvol}\,(\operatorname{div}v)(d^*A|_V,d^*A|_V)\\
&+\int\mathrm{dvol}\,(d^*A|_V,\iota_{\operatorname{grad}(\operatorname{div}v)}A|_V),
\end{aligned}
\]
and integration by parts cancels the first two terms, since $\int\mathrm{dvol}\,v(\cdot)=-\int\mathrm{dvol}\,(\operatorname{div}v)(\cdot)$. For the special conformal field, $\operatorname{div}v=-8(v_1,x)$ (Chapter 2) and $\operatorname{grad}(\operatorname{div}v)=-8v_1$, which gives the stated value. For translations and rotations $\operatorname{div}v=0$, and for $2E$, $\operatorname{div}v=8$ is constant.
\end{proof}

The finite counterterm of Theorem \ref{thm:c4-finite} is the Feynman-gauge longitudinal term, and Proposition \ref{prop:c4-K} shows that it breaks the special conformal symmetry. In Step 2 the components of the equivariant counterterms along the special conformal generators have to absorb this variation, or exhibit it as part of an obstruction.

\begin{remark}[Method of computation]
\label{rem:c4-method}
Each per-graph coefficient of this chapter is an exact rational number. It is obtained in an explicit realisation of $V_{10}$, $\Sigma$ and $\mathfrak g$ over $\mathbb Q(i)$, with the loop momentum integrated as a Gaussian (Lemma \ref{lem:c4-gauss}) and the $\varepsilon$-expansion taken on the graphs with one pinned edge (Proposition \ref{prop:c4-boundary}), in exact rational arithmetic; for the two-point graphs the polarisations and momenta are symbolic. For the cubic and quartic sets each multilinear form is determined by exact linear algebra on a basis of the $\mathrm{Spin}(4,\mathbb C)\times\mathrm{Spin}(6,\mathbb C)$-invariant multilinear forms of the given legs, including the parity-odd ones, from evaluations at integer polarisations and momenta drawn uniformly from a finite set of integers. A nonzero polynomial vanishes at such a point with probability at most its total degree divided by the size of the set (Schwartz--Zippel), so the sampled system has full rank with high probability; full rank is then established exactly, and the solution is exact. The colour factors are reduced exactly for a general simple $\mathfrak g$, and the resulting rational multiples of $\operatorname{tr}_{\mathfrak g}(\mathrm{ad}_X\mathrm{ad}_Y)/(X,Y)$ are the same in the explicit realisations $\mathfrak g=\mathfrak{su}(2)$ and $\mathfrak g=\mathfrak{su}(3)$. The sector $(c,A|_V)$ of the obstruction is also given by $Q$ of the renormalised two-point function (Remark \ref{rem:c4-selfenergy}).
\end{remark}
